% ===========================================================================
% Chapter 6 -- Deep Learning
% Rebuilt in the v3 didactic standard: every mechanism gets a plain-language
% model, a formal statement, a worked numeric example, an adjacent figure,
% a failure mode, and a retrieval prompt.
% ===========================================================================

\chapter{Deep Learning}
\label{ch:deep-learning}

\chapterlead{How do layered differentiable models learn representations at
scale?}{Represent}{Differentiate}{Optimise}{Scale}{deep-learning}

\begin{learningobjectives}
After working through this chapter you should be able to:
\begin{itemize}
  \item explain why a stack of linear layers collapses, and prove it in two
        lines;
  \item compute the output shape and parameter count of a dense, convolutional,
        or attention layer before writing any code;
  \item execute one backpropagation step by hand and verify that the loss
        decreases;
  \item calculate a scaled dot-product attention output from three tokens with
        pencil and paper;
  \item name the assumption that each regulariser enforces, and the failure it
        cannot prevent;
  \item choose between an RNN, a CNN, and a Transformer from the structure of
        the data rather than from fashion;
  \item estimate inference memory from a key--value cache and explain why
        parameter count alone is not a deployment metric;
  \item distinguish claims supported by primary evidence from claims that are
        folklore.
\end{itemize}
\end{learningobjectives}

\begin{plainlanguage}
A deep network is a long recipe made only of two repeating steps: multiply by
a matrix, then bend the result with a simple nonlinear function. Nothing in the
recipe is individually clever. What makes it powerful is that every number in
every matrix can be nudged, and that we have an efficient bookkeeping method
(backpropagation) for working out which way to nudge each one. Learning is a
very large number of very small, very well-informed corrections.
\end{plainlanguage}

\begin{engineernote}
Keep one habit from the first page to the last: \emph{track the shape}. Most
deep-learning bugs are not conceptual, they are a transposed matrix or a
forgotten batch dimension. If you can state the shape of the tensor entering
and leaving every block, you can debug almost anything in this chapter.
\end{engineernote}

% ===========================================================================
\section{From Biological Neurons to the Perceptron}
% ===========================================================================

\subsection{The Biological Starting Point}

Artificial neural networks borrow a metaphor, not a
blueprint\textsuperscript{\cite[][pp.~27--35]{Kriesel2007NeuralNetworks}}. A
biological neuron collects signals through \keyterm{dendrites}, integrates them
in the \keyterm{soma}, and---if the integrated membrane potential crosses a
threshold---emits an \keyterm{action potential} down its \keyterm{axon} to
\keyterm{synapses} that contact other neurons. The resting potential sits near
$-70\,\text{mV}$ and the firing threshold near
$-55\,\text{mV}$\textsuperscript{\cite[][pp.~29--31]{Kriesel2007NeuralNetworks}}.
Synaptic strengths change over time, which is the biological anchor for the
word \emph{learning}.

\begin{figure}[htbp]
\centering
\begin{tikzpicture}[
  >={Stealth[length=2mm]},
  soma/.style={circle,draw=AlmanachTeal,fill=AlmanachMist,minimum size=13mm,
    font=\sffamily\scriptsize,align=center},
  lbl/.style={font=\sffamily\scriptsize,color=AlmanachInk!75},
  sig/.style={->,thick,draw=AlmanachOchre}]
  \node[soma] (s) at (0,0) {soma};
  \foreach \a/\n in {150/1,180/2,210/3}{
    \draw[sig] ($(s)+(\a:20mm)$) -- ($(s)+(\a:7.5mm)$);
    \fill[AlmanachTeal!70] ($(s)+(\a:20mm)$) circle (0.7mm);}
  \node[lbl,left=21mm of s,align=center] {dendrites\\(inputs $x_i$)};
  \draw[very thick,draw=AlmanachTeal] (s.east) -- ++(26mm,0) coordinate (ax);
  \node[lbl,above] at ($(s.east)+(13mm,1mm)$) {axon};
  \foreach \d in {-4,0,4}{\draw[sig] (ax) -- ++(7mm,\d mm);}
  \node[lbl,right=8mm of ax,align=center] {synapses\\(outputs)};
  \draw[decorate,decoration={brace,amplitude=3pt,mirror},AlmanachInk!55]
    ($(s.south)+(-8mm,-3mm)$) -- ($(s.south)+(8mm,-3mm)$)
    node[midway,below=3pt,lbl,align=center]
    {integrate, then fire\\if $>-55$\,mV};
\end{tikzpicture}
\caption{The biological neuron reduced to the three operations an artificial
unit keeps: weighted collection of inputs, integration into a single scalar,
and a threshold that decides whether to pass a signal on. Everything else about
the cell is discarded by the model.}
\label{fig:dl-biological-neuron}
\end{figure}

\begin{v3aside}[The 100-step rule]
The brain performs complex recognition in roughly $500\,\text{ms}$, while
neurons fire at roughly $100\,\text{Hz}$. At most about 100 sequential
processing steps therefore fit into a recognition
event\textsuperscript{\cite[][pp.~29--31]{Kriesel2007NeuralNetworks}}. Whatever
the brain is doing, it cannot be a long serial program. It must be massively
parallel and shallow in sequence---which is exactly the shape of a layered
network evaluated on parallel hardware.
\end{v3aside}

\begin{cautionbox}[The metaphor has a hard limit]
Do not reason about artificial networks from neuroscience. Real neurons spike
in continuous time, learn through local chemistry, and do not compute gradients
of a global loss. The resemblance ends after \cref{fig:dl-biological-neuron}.
When this chapter says ``neuron'' it means a weighted sum followed by a scalar
function, nothing more.
\end{cautionbox}

\subsection{The Perceptron}

Rosenblatt's \keyterm{perceptron} is the smallest complete
unit\textsuperscript{\cite{rosenblattPerceptronProbabilisticModel1958},
\cite[][pp.~57--62]{Kriesel2007NeuralNetworks}}. It scores an input by a
weighted sum and thresholds the result:
\begin{equation}
  y =
  \begin{cases}
    1 & \text{if } \mathbf{w}^{\!\top}\mathbf{x} + b \geq 0,\\[2pt]
    0 & \text{otherwise.}
  \end{cases}
  \label{eq:dl-perceptron}
\end{equation}
\where{$\mathbf{x}$ is the input feature vector, $\mathbf{w}$ the learned
weight vector of the same length, and $b$ the bias that shifts the decision
threshold away from the origin.}

Geometrically, $\mathbf{w}^{\!\top}\mathbf{x} + b = 0$ is a
\keyterm{hyperplane}. The weight vector $\mathbf{w}$ is its normal direction and
$b$ shifts it away from the origin. The \keyterm{perceptron convergence
theorem} guarantees that the learning rule finds such a separating hyperplane in
a finite number of updates \emph{provided one exists}.

\begin{plainlanguage}
A perceptron is a bouncer with a clipboard. Each guest property (height, age,
shoe colour) gets a score weight, positive or negative. The bouncer adds up the
weighted scores, compares the total against a fixed cut-off, and says yes or
no. Training means adjusting the weights every time the bouncer gets one wrong.
\end{plainlanguage}

\subsection{XOR: The Failure That Started a Winter}

Minsky and Papert showed that a single perceptron cannot represent the
exclusive-or function\textsuperscript{\cite{minskyPerceptronsIntroductionComputational1969},
\cite[][pp.~63--65]{Kriesel2007NeuralNetworks}}. The reason is not subtle: XOR
is not linearly separable, and \cref{eq:dl-perceptron} can only carve the input
space with one straight cut. The result contributed to the first ``AI winter''.

\begin{workedexample}[Solving XOR with two hidden units, by hand]
The cure is one hidden layer with a nonlinearity. Take
$\mathrm{ReLU}(z)=\max(0,z)$ and fix these weights---nothing is learned here,
we are only proving that a solution \emph{exists}:
\[
  h_1 = \mathrm{ReLU}(x_1 + x_2 - 0),\qquad
  h_2 = \mathrm{ReLU}(x_1 + x_2 - 1),\qquad
  y   = h_1 - 2h_2 .
\]
Now evaluate all four cases:

\smallskip
\begin{tblr}{
  width=0.86\linewidth,
  colspec={Q[c]Q[c]Q[c]Q[c]Q[c]Q[c]},
  row{1}={font=\bfseries,bg=AlmanachMist},
  hlines,vlines}
$x_1$ & $x_2$ & $x_1{+}x_2$ & $h_1$ & $h_2$ & $y = h_1 - 2h_2$ \\
0 & 0 & 0 & 0 & 0 & $0$ \\
0 & 1 & 1 & 1 & 0 & $1$ \\
1 & 0 & 1 & 1 & 0 & $1$ \\
1 & 1 & 2 & 2 & 1 & $2-2 = 0$ \\
\end{tblr}
\smallskip

The output column is exactly XOR. Notice \emph{what the hidden layer did}: both
units read the same quantity $x_1+x_2$, but with different thresholds. Unit
$h_1$ fires whenever anything is on; unit $h_2$ fires only when \emph{both} are
on. Subtracting twice $h_2$ cancels the double-on case. The hidden layer did not
memorise four rows---it invented the feature ``how many inputs are on'', which
\emph{is} linearly separable.
\end{workedexample}

\begin{figure}[htbp]
\centering
\begin{tikzpicture}[scale=1.0,
  pt/.style={circle,draw=AlmanachInk,minimum size=4mm,inner sep=0pt,
    font=\scriptsize},
  cls0/.style={pt,fill=AlmanachMist},
  cls1/.style={pt,fill=AlmanachOchre!70},
  ax/.style={->,draw=AlmanachInk!60,>={Stealth[length=1.6mm]}},
  lbl/.style={font=\sffamily\scriptsize,color=AlmanachInk!75}]

  \begin{scope}
    \draw[ax] (-0.4,0) -- (2.5,0) node[right,lbl] {$x_1$};
    \draw[ax] (0,-0.4) -- (0,2.5) node[above,lbl] {$x_2$};
    \node[cls0] at (0,0) {}; \node[cls1] at (0,2) {};
    \node[cls1] at (2,0) {}; \node[cls0] at (2,2) {};
    \draw[thick,AlmanachRed,dashed] (-0.3,2.3) -- (2.4,-0.4);
    \draw[thick,AlmanachRed,dashed] (-0.3,1.2) -- (2.4,-1.5);
    \node[lbl,align=center] at (1.0,-1.1)
      {\textbf{Input space}\\no single line works};
  \end{scope}

  \begin{scope}[xshift=62mm]
    \draw[ax] (-0.4,0) -- (2.9,0) node[right,lbl] {$h_1$};
    \draw[ax] (0,-0.4) -- (0,2.5) node[above,lbl] {$h_2$};
    \node[cls0] at (0,0) {};
    \node[cls1] at (1,0) {};
    \node[cls0] at (2,1) {};
    \node[lbl] at (1.0,0.42) {$(0,1)$ and $(1,0)$ coincide};
    \draw[thick,AlmanachTeal] (-0.3,-0.55) -- (2.6,1.9);
    \node[lbl,align=center] at (1.2,-1.1)
      {\textbf{Hidden space}\\one line separates};
  \end{scope}
\end{tikzpicture}
\caption{Why depth helps. XOR is unsolvable by one straight cut in the input
space (left), but the hidden layer of the worked example re-coordinates the
four points so that a single line does separate them (right). Depth does not add
raw power to the classifier; it changes the space the classifier works in.}
\label{fig:dl-xor-transform}
\end{figure}

\begin{figure}[htbp]
\centering
\begin{tikzpicture}[
  node distance=9mm and 15mm,
  neuron/.style={circle,draw=AlmanachTeal,minimum size=7mm,font=\scriptsize,
    inner sep=1pt},
  inn/.style={neuron,fill=AlmanachMist},
  hid/.style={neuron,fill=AlmanachOchre!25},
  outn/.style={neuron,fill=AlmanachTeal!22},
  arr/.style={-{Stealth[length=1.5mm]},draw=AlmanachInk!45},
  lbl/.style={font=\sffamily\scriptsize,color=AlmanachInk!70}]
  \node[inn] (i1) {$x_1$};
  \node[inn,below=of i1] (i2) {$x_2$};
  \node[hid,right=of i1,yshift=-4.5mm] (h1) {$h_1$};
  \node[hid,below=of h1] (h2) {$h_2$};
  \node[outn,right=of h1,yshift=-4.5mm] (o) {$y$};
  \foreach \i in {i1,i2}\foreach \h in {h1,h2}{\draw[arr] (\i) -- (\h);}
  \draw[arr] (h1) -- (o) node[midway,above,lbl] {$+1$};
  \draw[arr] (h2) -- (o) node[midway,below,lbl] {$-2$};
  \node[lbl,below=2mm of i2] {input};
  \node[lbl,below=2mm of h2] {hidden (ReLU)};
  \node[lbl,below=11mm of o] {output};
  \node[lbl,above=2mm of h1] {bias $0$};
  \node[lbl,below left=1mm and -2mm of h2] {bias $-1$};
\end{tikzpicture}
\caption{The two-unit network that realises XOR, with the exact weights of the
worked example. Every edge from the input layer carries weight $+1$; the two
hidden units differ only in their bias, and the output layer subtracts twice the
second unit.}
\label{fig:dl-xor-network}
\end{figure}

\begin{checkpoint}
Change the output layer of the worked example to $y = h_1 - h_2$ and recompute
all four rows. Which logical function do you get now, and what does that tell
you about how much of the network's behaviour lives in the last layer?
\end{checkpoint}

% ===========================================================================
\section{Activation Functions and Why Nonlinearity Is Mandatory}
% ===========================================================================

\subsection{The Collapse Argument}

\begin{workedexample}[Two linear layers are one linear layer]
Suppose we drop the activation entirely. Layer one computes
$\mathbf{h} = W_1\mathbf{x} + \mathbf{b}_1$ and layer two computes
$\mathbf{y} = W_2\mathbf{h} + \mathbf{b}_2$. Substitute:
\[
  \mathbf{y}
  = W_2\left(W_1\mathbf{x} + \mathbf{b}_1\right) + \mathbf{b}_2
  = \underbrace{\left(W_2W_1\right)}_{=:W'}\mathbf{x}
  + \underbrace{\left(W_2\mathbf{b}_1 + \mathbf{b}_2\right)}_{=:\mathbf{b}'} .
\]
The composition is a single affine map $W'\mathbf{x}+\mathbf{b}'$. Stacking
$100$ such layers gives another single affine map. Depth without nonlinearity
buys \emph{nothing}: no extra expressiveness, only extra
multiplications\textsuperscript{\cite[][p.~17]{fleuretLittleBookDeep}}.
\end{workedexample}

This is the whole justification for activation functions. They are the only
reason depth is not a waste of electricity.

\subsection{The Standard Repertoire}

\begin{table}[htbp]
\centering
\caption{Activation functions and the engineering property that decides when
each is appropriate. The ``gradient risk'' column is the reason most modern
hidden layers use a rectifier family rather than a saturating sigmoid.}
\label{tab:dl-activations}
\begin{tblr}{
  width=\textwidth,
  colspec={Q[l,0.15\textwidth] Q[l,0.24\textwidth] Q[c,0.15\textwidth] X[l]},
  row{1}={font=\bfseries,bg=AlmanachMist},
  hlines,vlines}
Function & Definition & Range & Gradient risk and typical use \\
Sigmoid & $\sigma(z)=\dfrac{1}{1+e^{-z}}$ & $(0,1)$ &
  Saturates both sides; $\sigma'\le 0.25$, so gradients shrink by
  ${\ge}4\times$ per layer. Use only for a binary output. \\
Tanh & $\tanh(z)$ & $(-1,1)$ &
  Zero-centred, still saturates. Common in classic RNN cells. \\
ReLU & $\max(0,z)$ & $[0,\infty)$ &
  Gradient is exactly $0$ or $1$: no shrinkage on the active side. Default
  hidden activation\textsuperscript{\cite{glorotUnderstandingDifficultyTraining}}. \\
Leaky ReLU & $\max(\alpha z, z)$, $\alpha\!\approx\!0.01$ &
  $(-\infty,\infty)$ &
  Keeps a small negative gradient so units cannot die permanently. \\
GELU & $z\,\Phi(z)$ & ${\approx}(-0.17,\infty)$ &
  Smooth near the origin; standard inside Transformer feed-forward blocks. \\
Softmax & $\dfrac{e^{z_i}}{\sum_j e^{z_j}}$ & $(0,1)$, sums to $1$ &
  Output layer for multi-class problems. Not an activation for hidden layers. \\
\end{tblr}
\end{table}

\begin{figure}[htbp]
\centering
\begin{tikzpicture}
\begin{axis}[almanach axis,
  height=44mm,
  domain=-4:4,
  samples=161,
  ymin=-1.3,ymax=2.6,
  xlabel={pre-activation $z$},
  ylabel={output},
  legend pos=north west,
  legend columns=2]
  \addplot[AlmanachTeal,thick]   {1/(1+exp(-x))};   \addlegendentry{sigmoid}
  \addplot[AlmanachOchre,thick]  {tanh(x)};         \addlegendentry{tanh}
  \addplot[AlmanachRed,thick]    {max(0,x)};        \addlegendentry{ReLU}
  \addplot[AlmanachInk,thick,dashed]
    {0.5*x*(1+tanh(0.7978845608*(x+0.044715*x^3)))}; \addlegendentry{GELU}
\end{axis}
\end{tikzpicture}
\hfill
\begin{tikzpicture}
\begin{axis}[almanach axis,
  height=44mm,
  domain=-4:4,
  samples=161,
  ymin=-0.1,ymax=1.15,
  xlabel={pre-activation $z$},
  ylabel={derivative},
  legend pos=north west]
  \addplot[AlmanachTeal,thick]  {exp(-x)/(1+exp(-x))^2};
    \addlegendentry{$\sigma'$}
  \addplot[AlmanachOchre,thick] {1-(tanh(x))^2}; \addlegendentry{$\tanh'$}
  \addplot[AlmanachRed,thick]   {x>0};           \addlegendentry{ReLU$'$}
\end{axis}
\end{tikzpicture}
\caption{The functions (left) look broadly similar; their derivatives (right)
explain everything. The sigmoid derivative never exceeds $0.25$ and collapses
toward zero away from the origin, so a deep stack multiplies many small numbers
together. The ReLU derivative is $1$ wherever the unit is active, which is why
rectifiers made deep stacks trainable.}
\label{fig:dl-activation-curves}
\end{figure}

\begin{workedexample}[Why saturation kills depth: multiply the numbers]
Backpropagation multiplies one derivative per layer (\cref{sec:dl-backprop}).
Suppose every layer sits in a mildly saturated sigmoid regime with
$\sigma' \approx 0.2$. After $10$ layers the gradient reaching layer one is
scaled by
\[
  0.2^{10} = 1.024\times10^{-7},
\]
and after $20$ layers by $1.05\times10^{-14}$. In \texttt{float32}, whose
resolution near $1$ is about $1.2\times10^{-7}$, the early layers have stopped
receiving a usable signal. Repeat with ReLU units that are active,
$\left(1\right)^{20} = 1$: the signal arrives undiminished. This single
calculation is the practical content of the phrase \emph{vanishing gradient}.
\end{workedexample}

\begin{engineernote}
Default to ReLU or GELU in hidden layers, sigmoid only on a binary output,
softmax only on a multi-class output. Choosing an exotic activation is almost
never the highest-value change available to you; fixing the learning rate,
the data, or the normalisation usually is.
\end{engineernote}

\begin{cautionbox}[Dying ReLU]
If a unit's pre-activation is negative for \emph{every} training example, its
gradient is exactly zero forever and the unit is dead weight---it can never
recover, because the gradient that would move it is itself zero. A too-large
learning rate is the usual cause: one enormous update pushes the bias far
negative. Symptoms: a fraction of units with permanently zero activation, and a
loss that plateaus above the achievable value. Remedies: lower the learning
rate, use Leaky ReLU or GELU, and check He initialisation
(\cref{sec:dl-init}).
\end{cautionbox}

\subsection{Softmax and Cross-Entropy}

For $K$ classes, the \keyterm{softmax} turns raw scores (\keyterm{logits}) into
a distribution\textsuperscript{\cite[][p.~18]{fleuretLittleBookDeep}}:
\begin{equation}
  p_i = \mathrm{softmax}(z)_i = \frac{e^{z_i}}{\sum_{j=1}^{K} e^{z_j}},
  \qquad p_i > 0,\quad \sum_{i=1}^{K} p_i = 1 .
  \label{eq:dl-softmax}
\end{equation}
Paired with \keyterm{cross-entropy} loss $\mathcal{L} = -\log p_{c}$ for the
true class $c$, the gradient with respect to the logits is remarkably clean:
\begin{equation}
  \frac{\partial \mathcal{L}}{\partial z_i} = p_i - y_i ,
  \label{eq:dl-softmax-grad}
\end{equation}
where $y$ is the one-hot target.

\begin{workedexample}[Softmax and cross-entropy with real numbers]
Let the logits be $z = (2.0,\; 1.0,\; 0.1)$ and let the true class be the first
one, $y = (1,0,0)$.

\emph{Step 1---exponentiate.}
$e^{2.0}=7.389$, $e^{1.0}=2.718$, $e^{0.1}=1.105$.

\emph{Step 2---normalise.} The sum is $7.389+2.718+1.105 = 11.212$, so
\[
  p = \left(\tfrac{7.389}{11.212},\ \tfrac{2.718}{11.212},\
  \tfrac{1.105}{11.212}\right) = (0.659,\ 0.242,\ 0.099).
\]
Check: $0.659+0.242+0.099 = 1.000$. \checkmark

\emph{Step 3---loss.} $\mathcal{L} = -\log(0.659) = 0.417$ nats.

\emph{Step 4---gradient.} By \cref{eq:dl-softmax-grad},
\[
  \frac{\partial \mathcal{L}}{\partial z}
  = p - y = (0.659-1,\ 0.242,\ 0.099)
  = (-0.341,\ +0.242,\ +0.099).
\]
Check: the components sum to $0.000$. \checkmark\ They must, because softmax
outputs are constrained to sum to one---pushing one logit up necessarily pushes
the others down. \emph{Read the signs}: the gradient is negative for the true
class (increase that logit) and positive for the wrong ones (decrease them).
That is the entire learning signal of a classifier.
\end{workedexample}

\begin{cautionbox}[Softmax outputs are not calibrated probabilities]
A softmax output of $0.659$ does \emph{not} mean the model is right $65.9\,\%$
of the time on inputs like this one. Modern networks are typically
overconfident. Calibration is a separate empirical property that must be
measured on held-out data (see \cref{ch:evaluation-mlops}) and, if needed,
corrected---for example by temperature scaling. Never quote a raw softmax value
as a probability in a report without a calibration check.
\end{cautionbox}

\begin{checkpoint}
Add the same constant $c=10$ to all three logits in the worked example and
recompute $p$. Why is the result unchanged, and why do numerical libraries
subtract $\max_i z_i$ before exponentiating?
\end{checkpoint}

% ===========================================================================
\section{Feed-Forward Networks, Shapes, and Capacity}
% ===========================================================================

\subsection{The Dense Layer}

A \keyterm{feed-forward} (fully connected, dense) layer applies one affine map
followed by one activation\textsuperscript{\cite[][pp.~15--17]{fleuretLittleBookDeep}}:
\begin{equation}
  \begin{gathered}
    \mathbf{h} = \phi\!\left(W\mathbf{x} + \mathbf{b}\right),\\[2pt]
    W \in \mathbb{R}^{n_{\text{out}}\times n_{\text{in}}},\qquad
    \mathbf{b}\in\mathbb{R}^{n_{\text{out}}} .
  \end{gathered}
  \label{eq:dl-dense}
\end{equation}
\where{$\mathbf{x}$ is the layer's input, $W$ the weight matrix, $\mathbf{b}$
the bias vector, and $\phi$ the activation applied element-wise;
$n_{\text{in}}$ and $n_{\text{out}}$ are the input and output widths.}
Information flows strictly forward; there are no cycles.

\begin{workedexample}[Counting parameters before writing code]
A layer mapping a flattened $28\times28$ greyscale image to $128$ hidden units
has $n_{\text{in}} = 784$ and $n_{\text{out}} = 128$:
\[
  \underbrace{784\times128}_{\text{weights } W} +
  \underbrace{128}_{\text{biases } \mathbf{b}} = 100{,}352 + 128 = 100{,}480
  \text{ parameters.}
\]
At $4$ bytes each in \texttt{float32} that is $401{,}920$ bytes, about
$0.4\,\text{MB}$---for \emph{one} layer. Now add the optimiser: Adam stores two
extra moment estimates per parameter, so training memory is roughly
$3\times$ the parameter memory before activations are counted at all. This is
why ``how big is the model'' and ``how much memory does training need'' are
different questions with answers that differ by an order of magnitude.
\end{workedexample}

\begin{figure}[htbp]
\centering
\begin{tikzpicture}[
  node distance=6mm,
  blk/.style={draw=AlmanachTeal,fill=AlmanachMist,rounded corners=1mm,
    minimum height=8mm,minimum width=21mm,font=\sffamily\scriptsize,
    align=center},
  act/.style={blk,fill=AlmanachOchre!20,draw=AlmanachOchre,minimum width=13mm},
  shp/.style={font=\ttfamily\scriptsize,color=AlmanachTeal!75!black},
  arr/.style={-{Stealth[length=1.8mm]},thick,draw=AlmanachInk!55}]
  \node[blk] (x) {input};
  \node[blk,right=of x] (l1) {dense\\$784\!\to\!128$};
  \node[act,right=of l1] (a1) {ReLU};
  \node[blk,right=of a1] (l2) {dense\\$128\!\to\!10$};
  \node[act,right=of l2] (a2) {softmax};
  \foreach \a/\b in {x/l1,l1/a1,a1/l2,l2/a2}{\draw[arr] (\a) -- (\b);}
  \node[shp,below=1.5mm of x]  {(B, 784)};
  \node[shp,below=1.5mm of l1] {(B, 128)};
  \node[shp,below=1.5mm of a1] {(B, 128)};
  \node[shp,below=1.5mm of l2] {(B, 10)};
  \node[shp,below=1.5mm of a2] {(B, 10)};
  \node[font=\sffamily\scriptsize,color=AlmanachInk!70,above=1.5mm of l1]
    {100{,}480 params};
  \node[font=\sffamily\scriptsize,color=AlmanachInk!70,above=1.5mm of l2]
    {1{,}290 params};
\end{tikzpicture}
\caption{A shape-tracking diagram for a minimal classifier, with the batch
dimension $B$ carried explicitly. Annotating every arrow with its tensor shape
turns most debugging into arithmetic: if two adjacent shapes do not match, the
bug is located before the code is ever run.}
\label{fig:dl-shape-tracking}
\end{figure}

\subsection{Universal Approximation, and What It Does Not Promise}

Cybenko and Hornik proved that a feed-forward network with a single hidden
layer and enough units can approximate any continuous function on a compact set
to arbitrary accuracy\textsuperscript{\cite{cybenkoApproximationSuperpositionsSigmoidal1989,
hornikApproximationCapabilitiesMultilayer1991}}.

\begin{cautionbox}[The most over-quoted theorem in the field]
Universal approximation is an \emph{existence} statement about representation.
It does not say that:
\begin{itemize}
  \item the required width is practical---it may grow exponentially with input
        dimension;
  \item gradient descent will \emph{find} the approximating weights;
  \item the found function will generalise beyond the training data.
\end{itemize}
Depth is the practical answer: for many function classes, deep networks are
exponentially more parameter-efficient than wide shallow
ones\textsuperscript{\cite[][pp.~8--12]{robertsPrinciplesDeepLearning2022}}.
That is the actual motivation for the word \emph{deep}. Anyone who cites this
theorem as evidence that a network will work on their problem has misread it.
\end{cautionbox}

\begin{checkpoint}
Your colleague proposes a single hidden layer with $10^6$ units instead of ten
layers of $10^3$ units, citing universal approximation. Give two concrete
engineering objections and one situation in which they would nevertheless be
right.
\end{checkpoint}

% ===========================================================================
\section{Training: Loss, Backpropagation, and Optimisation}
\label{sec:dl-backprop}
% ===========================================================================

\subsection{The Forward Pass and the Loss}

In the \keyterm{forward pass}, data flows through the layers and the final
prediction $\hat y$ is compared to the target $y$ by a \keyterm{loss function}
$\mathcal{L}(y,\hat y)$\textsuperscript{\cite[][pp.~20--22]{fleuretLittleBookDeep}}.
The loss is the only place where the engineer states what ``good'' means, and
it is chosen, not discovered.

\begin{table}[htbp]
\centering
\caption{Standard losses and the assumption each one encodes. Choosing a loss
is choosing which errors you are willing to tolerate.}
\label{tab:dl-losses}
\begin{tblr}{
  width=\textwidth,
  colspec={Q[l,0.17\textwidth] Q[l,0.29\textwidth] X[l]},
  row{1}={font=\bfseries,bg=AlmanachMist},
  hlines,vlines}
Loss & Definition & What it assumes, and its failure mode \\
MSE & $\frac1N\sum_i (y_i-\hat y_i)^2$ &
  Gaussian noise, symmetric cost. Quadratic penalty means a single outlier can
  dominate the gradient. \\
MAE & $\frac1N\sum_i \lvert y_i-\hat y_i\rvert$ &
  Laplacian noise; robust to outliers, but the gradient has constant magnitude
  and does not shrink near the optimum. \\
Cross-entropy & $-\sum_i y_i\log\hat y_i$ &
  The target is a distribution. Unbounded when the model is confidently wrong,
  which is what makes it train well. \\
Binary CE & $-\left[y\log\hat y + (1{-}y)\log(1{-}\hat y)\right]$ &
  Two outcomes. Silently misleading under strong class imbalance---report
  precision, recall, and a baseline alongside it. \\
\end{tblr}
\end{table}

\subsection{Backpropagation Is the Chain Rule with Bookkeeping}

\keyterm{Backpropagation}\textsuperscript{%
\cite{rumelhartLearningRepresentationsBackpropagating1986},
\cite[][pp.~1--4]{lecunTheoreticalFrameworkBack}} computes
$\partial\mathcal{L}/\partial\theta$ for every parameter by applying the chain
rule backwards from the output:
\begin{equation}
  \frac{\partial\mathcal{L}}{\partial W_l}
  = \frac{\partial\mathcal{L}}{\partial a_L}\,
    \frac{\partial a_L}{\partial a_{L-1}}\cdots
    \frac{\partial a_{l+1}}{\partial a_l}\,
    \frac{\partial a_l}{\partial W_l}.
  \label{eq:dl-backprop}
\end{equation}
\where{$\mathcal{L}$ is the loss, $W_l$ the weights of layer $l$, $a_l$ the
activations leaving layer $l$, and $L$ the index of the final layer. Each
factor is one layer's local derivative, and the product runs backwards from the
loss to the layer being updated.}
The algorithmic insight is that each intermediate factor is computed
\emph{once} and reused by every earlier layer. That reuse is what makes the cost
of the backward pass a small constant multiple of the forward pass rather than
proportional to the number of parameters.

\begin{plainlanguage}
Imagine a long assembly line that produces a defective part. To assign
responsibility you walk backwards from the inspector to the first station,
asking at each step: ``given how much the next station's output was off, how
much of that came from my settings?'' Backpropagation is exactly this walk, and
its efficiency comes from the fact that each station passes its answer to the
one before it instead of everyone recomputing the whole line.
\end{plainlanguage}

\begin{workedexample}[One complete backpropagation step, by hand]
Take the smallest network that still has a hidden layer:
\[
  z_1 = w_1x + b_1,\quad a_1 = \mathrm{ReLU}(z_1),\quad
  \hat y = z_2 = w_2a_1 + b_2,\quad
  \mathcal{L} = \tfrac12(\hat y - y)^2 .
\]
Initialise $w_1=0.5$, $b_1=0$, $w_2=-1$, $b_2=1$, and take the training pair
$x=2$, $y=1$, with learning rate $\eta = 0.1$.

\emph{Forward.}
$z_1 = 0.5\cdot2 + 0 = 1.0$;
$a_1 = \mathrm{ReLU}(1.0) = 1.0$;
$\hat y = -1\cdot1.0 + 1 = 0.0$;
$\mathcal{L} = \tfrac12(0-1)^2 = \mathbf{0.5}$.

\emph{Backward.} Start at the loss and walk left, writing
$\delta_u := \partial\mathcal{L}/\partial u$ to keep the lines short:

\smallskip
\begin{tblr}{
  width=0.95\linewidth,
  colspec={Q[l,0.13\linewidth] X[l] Q[r,0.13\linewidth]},
  row{1}={font=\bfseries,bg=AlmanachMist},
  hlines,vlines}
Quantity & Chain-rule expression & Value \\
$\delta_{\hat y}$ & $\hat y - y$ & $-1$ \\
$\delta_{w_2}$ & $\delta_{\hat y}\cdot a_1 = (-1)(1.0)$ & $-1$ \\
$\delta_{b_2}$ & $\delta_{\hat y}$ & $-1$ \\
$\delta_{a_1}$ & $\delta_{\hat y}\cdot w_2 = (-1)(-1)$ & $+1$ \\
$\delta_{z_1}$ & $\delta_{a_1}\cdot \mathrm{ReLU}'(z_1) = (1)(1)$ & $+1$ \\
$\delta_{w_1}$ & $\delta_{z_1}\cdot x = (1)(2)$ & $+2$ \\
$\delta_{b_1}$ & $\delta_{z_1}$ & $+1$ \\
\end{tblr}
\smallskip

\emph{Update} with $\theta \leftarrow \theta - \eta\,\partial\mathcal{L}/
\partial\theta$:
\[
  w_1 = 0.5 - 0.1(2) = 0.3,\quad b_1 = 0 - 0.1(1) = -0.1,\quad
  w_2 = -1 + 0.1 = -0.9,\quad b_2 = 1 + 0.1 = 1.1 .
\]

\emph{Verify.} Run the forward pass again:
$z_1 = 0.3\cdot2 - 0.1 = 0.5$, $a_1 = 0.5$,
$\hat y = -0.9\cdot0.5 + 1.1 = 0.65$, so
\[
  \mathcal{L}_{\text{new}} = \tfrac12(0.65-1)^2 = \tfrac12(0.1225)
  = \mathbf{0.061}.
\]
The loss fell from $0.5$ to $0.061$ in one step. \checkmark\ Notice that
$\partial\mathcal{L}/\partial w_1 = 2$ is the largest gradient in the network
purely because $x=2$ is the largest activation feeding it---gradients scale with
the activations they multiply, which is the reason input normalisation matters
so much.
\end{workedexample}

\begin{figure}[htbp]
\centering
\begin{tikzpicture}[
  node distance=7mm and 13mm,
  op/.style={draw=AlmanachTeal,fill=AlmanachMist,rounded corners=1mm,
    minimum height=7mm,minimum width=15mm,font=\sffamily\scriptsize},
  val/.style={font=\scriptsize,color=AlmanachTeal!75!black},
  grad/.style={font=\scriptsize,color=AlmanachRed},
  fw/.style={-{Stealth[length=1.8mm]},thick,draw=AlmanachTeal},
  bw/.style={-{Stealth[length=1.8mm]},thick,draw=AlmanachRed,dashed}]
  \node[op] (x)  {$x$};
  \node[op,right=of x]  (z1) {$z_1$};
  \node[op,right=of z1] (a1) {$a_1$};
  \node[op,right=of a1] (y)  {$\hat y$};
  \node[op,right=of y]  (L)  {$\mathcal{L}$};
  \foreach \a/\b in {x/z1,z1/a1,a1/y,y/L}{\draw[fw] (\a) -- (\b);}
  \foreach \a/\b in {L/y,y/a1,a1/z1,z1/x}{
    \draw[bw] ([yshift=-3.6mm]\a.south) -- ([yshift=-3.6mm]\b.south);}
  \node[val,above=1mm of x]  {$2.0$};
  \node[val,above=1mm of z1] {$1.0$};
  \node[val,above=1mm of a1] {$1.0$};
  \node[val,above=1mm of y]  {$0.0$};
  \node[val,above=1mm of L]  {$0.5$};
  \node[grad,below=5mm of z1] {$1$};
  \node[grad,below=5mm of a1] {$1$};
  \node[grad,below=5mm of y]  {$-1$};
  \node[font=\sffamily\scriptsize,color=AlmanachTeal,above=6mm of a1]
    {forward: values};
  \node[font=\sffamily\scriptsize,color=AlmanachRed,below=10mm of a1]
    {backward: gradients};
\end{tikzpicture}
\caption{The computation graph of the worked example, carrying forward values
above each node and backward gradients below. Every deep-learning framework
builds exactly this graph at run time; ``automatic differentiation'' is the
mechanical replay of the lower row.}
\label{fig:dl-computation-graph}
\end{figure}

\subsection{Optimisers}

Backpropagation supplies gradients; an \keyterm{optimiser} decides what to do
with them\textsuperscript{\cite[][pp.~22--25]{fleuretLittleBookDeep}}.

\begin{table}[htbp]
\centering
\caption{Optimisers as successive fixes to one problem: raw SGD steps are noisy
and badly scaled. Each row adds a correction and a new hyperparameter.}
\label{tab:dl-optimisers}
\begin{tblr}{
  width=\textwidth,
  colspec={Q[l,0.17\textwidth] Q[l,0.30\textwidth] X[l]},
  row{1}={font=\bfseries,bg=AlmanachMist},
  hlines,vlines}
Optimiser & Core update & What it fixes, and what it costs \\
SGD & $\theta \leftarrow \theta - \eta\nabla\mathcal{L}$ &
  Nothing yet. Cheap, unbiased, but noisy and sensitive to feature scaling. \\
SGD + momentum & $v \leftarrow \beta v + \nabla\mathcal{L}$;
  $\theta \leftarrow \theta - \eta v$ &
  Damps oscillation across narrow valleys. Adds $\beta$ and one extra buffer
  per parameter. \\
RMSprop & Divide by a running RMS of recent gradients &
  Per-parameter scaling, so rarely updated weights still move. \\
Adam & Momentum $+$ RMSprop with bias
  correction\textsuperscript{\cite{kingmaAdamMethodStochastic2017}} &
  Robust default. Costs \emph{two} extra buffers per parameter---triple the
  optimiser memory. \\
AdamW & Adam with weight decay applied
  separately\textsuperscript{\cite{loshchilovDecoupledWeightDecay2019}} &
  Fixes the fact that L2 inside Adam is \emph{not} equivalent to weight decay.
  Standard for Transformers. \\
\end{tblr}
\end{table}

\begin{engineernote}
Use AdamW unless you have a measured reason not to. The historically important
detail is the one Loshchilov and Hutter
identified\textsuperscript{\cite{loshchilovDecoupledWeightDecay2019}}: adding an
L2 term to the loss and calling it ``weight decay'' gives different behaviour
under Adam, because the adaptive denominator rescales the penalty too. If you
inherit code that sets \texttt{weight\_decay} on a plain \texttt{Adam}, that is
a real bug, not a stylistic difference.
\end{engineernote}

\subsection{Mini-Batches, Learning Rate, and Schedules}

Gradients are computed on \keyterm{mini-batches}, typically $32$ to $512$
examples\textsuperscript{\cite[][p.~23]{fleuretLittleBookDeep}}. This trades the
low noise and high cost of full-batch descent against the high noise and low
cost of single-sample updates. The residual noise is not purely a nuisance: it
helps escape saddle points and sharp minima.

The \keyterm{learning rate} $\eta$ is the single most consequential
hyperparameter. Too large and the loss diverges or oscillates; too small and
training stalls in an apparently converged plateau.

\begin{figure}[htbp]
\centering
\begin{tikzpicture}
\begin{axis}[almanach axis,
  height=44mm,
  domain=0:100,
  samples=201,
  xlabel={training step},
  ylabel={learning rate $\eta$},
  ymin=0,ymax=1.15,
  legend pos=north east]
  \addplot[AlmanachTeal,thick]
    {x<10 ? x/10 : 0.5*(1+cos(deg(pi*(x-10)/90)))};
    \addlegendentry{warmup + cosine}
  \addplot[AlmanachOchre,thick,dashed]
    {x<33 ? 1 : (x<66 ? 0.3 : 0.09)};
    \addlegendentry{step decay}
  \addplot[AlmanachRed,thick,dotted] {1};
    \addlegendentry{constant}
\end{axis}
\end{tikzpicture}
\caption{Three learning-rate schedules. The warmup ramp exists because the very
first updates act on randomly initialised weights where gradient estimates are
unreliable; the cosine decay then anneals the step size so late training can
settle into a minimum instead of bouncing around it.}
\label{fig:dl-lr-schedules}
\end{figure}

\begin{cautionbox}[Diagnose the learning rate before anything else]
When a network will not train, the learning rate is the first suspect, not the
architecture. Read the loss curve:
\begin{itemize}
  \item loss becomes \texttt{NaN} or rises steeply $\rightarrow$ $\eta$ far too
        large;
  \item loss falls then oscillates on a plateau $\rightarrow$ $\eta$ slightly
        too large, or the schedule decays too slowly;
  \item loss falls smoothly but far too slowly $\rightarrow$ $\eta$ too small;
  \item training loss will not fall even on a \emph{single batch}
        $\rightarrow$ the bug is not the learning rate. It is the data
        pipeline, the loss, or the label alignment.
\end{itemize}
That last test---can the model overfit one batch to near-zero loss?---is the
cheapest and most informative check in deep learning. Run it before every
serious experiment.
\end{cautionbox}

\begin{checkpoint}
You double the batch size from $64$ to $128$. State what happens to (a) the
variance of the gradient estimate, (b) the number of updates per epoch, and
(c) the learning rate you should now use. Justify (c) from (a).
\end{checkpoint}

% ===========================================================================
\section{Initialisation, Normalisation, and Residual Paths}
\label{sec:dl-init}
% ===========================================================================

\subsection{Why Initialisation Is Not Arbitrary}

Initialising all weights to zero creates a \keyterm{symmetry
problem}\textsuperscript{\cite[][pp.~8--10]{robertsPrinciplesDeepLearning2022}}:
every unit in a layer computes the same output, receives the same gradient, and
stays identical forever. The layer behaves as if it had one unit. Random
initialisation breaks the symmetry, but the \emph{scale} of the randomness
decides whether the signal survives depth.

\keyterm{Xavier/Glorot} initialisation targets sigmoid and tanh
layers\textsuperscript{\cite{glorotUnderstandingDifficultyTraining}}, and
\keyterm{He} initialisation corrects it for
rectifiers\textsuperscript{\cite{heDelvingDeepRectifiers2015}}:
\begin{equation}
  W_{\text{Xavier}} \sim
  \mathcal{N}\!\left(0,\ \frac{2}{n_{\text{in}}+n_{\text{out}}}\right),
  \qquad
  W_{\text{He}} \sim
  \mathcal{N}\!\left(0,\ \frac{2}{n_{\text{in}}}\right).
  \label{eq:dl-init}
\end{equation}

\begin{workedexample}[Where the factor of two comes from]
Consider one unit with $n_{\text{in}}$ inputs, weights drawn independently with
variance $\sigma_w^2$, and inputs with variance $\sigma_x^2$. For the
pre-activation $z=\sum_{i=1}^{n_{\text{in}}} w_ix_i$,
\[
  \operatorname{Var}(z) = n_{\text{in}}\,\sigma_w^2\,\sigma_x^2 .
\]
To keep $\operatorname{Var}(z) = \sigma_x^2$---signal neither growing nor
shrinking with depth---we need $\sigma_w^2 = 1/n_{\text{in}}$.

Now add ReLU. It zeroes the negative half, which for a symmetric distribution
removes about half the variance. Compensating requires doubling:
$\sigma_w^2 = 2/n_{\text{in}}$, which is exactly He initialisation.

\emph{Numerically}: with $n_{\text{in}}=256$ and unit-variance inputs, He gives
$\sigma_w = \sqrt{2/256} = 0.088$. If you instead used $\sigma_w = 1$, the
pre-activation variance would be $256$, so activations would be around
$\pm16$---deep into saturation for any sigmoid and numerically hostile for any
activation. Over $10$ layers the blow-up factor is $256^{5} \approx
1.1\times10^{12}$.
\end{workedexample}

\subsection{Normalisation Layers}

\keyterm{Batch normalisation} normalises each feature across the
mini-batch\textsuperscript{\cite{ioffeBatchNormalizationAccelerating2015}}, then
restores expressiveness with learned scale $\gamma$ and shift $\beta$:
\begin{equation}
  \hat x_i = \frac{x_i-\mu_{\mathcal{B}}}
                  {\sqrt{\sigma^2_{\mathcal{B}}+\varepsilon}},
  \qquad
  y_i = \gamma\hat x_i + \beta .
  \label{eq:dl-batchnorm}
\end{equation}
\where{$x_i$ is one activation, $\mu_{\mathcal{B}}$ and
$\sigma^2_{\mathcal{B}}$ the mean and variance over the current mini-batch
$\mathcal{B}$, and $\varepsilon$ a small constant preventing division by zero.
$\gamma$ and $\beta$ are learned, so the layer can undo the normalisation if
that is what helps.}
\keyterm{Layer normalisation} performs the same computation but across the
\emph{features of a single example} rather than across the
batch\textsuperscript{\cite{baLayerNormalization2016}}. That difference matters
enormously: layer norm is independent of batch size and of other examples, which
is why Transformers use it and why it works in autoregressive generation with
batch size one.

\begin{workedexample}[Layer normalisation on four numbers]
Let one token's feature vector be $x = (2,\,4,\,6,\,8)$.

\emph{Mean}: $\mu = (2+4+6+8)/4 = 5$.

\emph{Variance}: $\sigma^2 = \frac{(-3)^2+(-1)^2+1^2+3^2}{4}
= \frac{9+1+1+9}{4} = 5$, so $\sigma = 2.236$.

\emph{Normalise}: $\hat x = \left(\frac{-3}{2.236},\frac{-1}{2.236},
\frac{1}{2.236},\frac{3}{2.236}\right) = (-1.342,\,-0.447,\,0.447,\,1.342)$.

Check: the mean is $0$ and the variance is $1$. \checkmark\ Every token now
enters the next layer at a comparable scale regardless of how large its raw
activations were, which is precisely the property that keeps a $30$-layer stack
numerically stable.
\end{workedexample}

\begin{cautionbox}[Batch norm has a train/test asymmetry]
During training, batch norm uses the statistics of the current mini-batch. At
inference it must use \emph{running} estimates accumulated during training. Two
classic failures follow: forgetting to switch the module into evaluation mode
(so predictions depend on which other examples happen to share the batch---a
genuine information leak between test samples), and using very small batches,
where the batch statistics are too noisy to be useful. If either applies,
prefer layer norm or group norm.
\end{cautionbox}

\subsection{Residual Connections}

A \keyterm{residual} block computes $y = x + F(x)$ rather than $y =
F(x)$\textsuperscript{\cite{heDeepResidualLearning2016}}. Differentiating shows
why this single change made networks with hundreds of layers trainable:
\begin{equation}
  \frac{\partial y}{\partial x} = I + \frac{\partial F}{\partial x}.
  \label{eq:dl-residual}
\end{equation}
\where{$F$ is the learned branch of the block, $x$ its input, and $I$ the
identity matrix. The $I$ term is the point: it guarantees a gradient path of
magnitude one regardless of what $F$ does.}
The identity term guarantees a gradient path of magnitude one straight through
the block. Even if the learned branch $F$ contributes nothing---at
initialisation it very nearly does not---the gradient still reaches the layers
below undiminished.

\begin{figure}[htbp]
\centering
\begin{tikzpicture}[
  node distance=6mm,
  blk/.style={draw=AlmanachTeal,fill=AlmanachMist,rounded corners=1mm,
    minimum height=7mm,minimum width=20mm,font=\sffamily\scriptsize,
    align=center},
  plus/.style={circle,draw=AlmanachOchre,fill=white,inner sep=1pt,
    font=\scriptsize},
  arr/.style={-{Stealth[length=1.8mm]},thick,draw=AlmanachInk!55},
  hi/.style={-{Stealth[length=1.8mm]},thick,draw=AlmanachOchre},
  lbl/.style={font=\sffamily\scriptsize,color=AlmanachInk!70}]

  \node[lbl] (t1) at (0,2.3) {\textbf{Plain block}};
  \node[blk] (px) at (0,1.5) {$x$};
  \node[blk,below=of px] (pf) {$F(\cdot)$};
  \node[blk,below=of pf] (py) {$y=F(x)$};
  \draw[arr] (px) -- (pf); \draw[arr] (pf) -- (py);
  \node[lbl,right=1mm of pf,align=left,text width=17mm]
    {$\partial y/\partial x=\partial F/\partial x$};

  \node[lbl] (t2) at (6.3,2.3) {\textbf{Residual block}};
  \node[blk] (rx) at (6.3,1.5) {$x$};
  \node[blk,below=of rx] (rf) {$F(\cdot)$};
  \node[plus,below=of rf] (rp) {$+$};
  \node[blk,below=of rp] (ry) {$y=x+F(x)$};
  \draw[arr] (rx) -- (rf); \draw[arr] (rf) -- (rp);
  \draw[arr] (rp) -- (ry);
  \draw[hi] (rx.west) -- ++(-9mm,0) |- (rp.west);
  \node[lbl,left=1mm of rf,align=right,text width=13mm,color=AlmanachOchre]
    {identity\\shortcut};
  \node[lbl,right=1mm of rp,align=left,text width=18mm]
    {$\partial y/\partial x=I+\partial F/\partial x$};
\end{tikzpicture}
\caption{The residual shortcut in one picture. In the plain block the gradient
must survive multiplication by every Jacobian on the way down; in the residual
block the identity term adds a lossless path, so depth no longer attenuates the
learning signal.}
\label{fig:dl-residual}
\end{figure}

\begin{checkpoint}
Twenty stacked plain blocks each have $\lVert\partial F/\partial x\rVert
\approx 0.7$. Estimate the gradient magnitude reaching the first block. Repeat
for twenty residual blocks using \cref{eq:dl-residual}. Which number explains
the depth records set after 2015?
\end{checkpoint}

% ===========================================================================
\section{Regularisation and Generalisation}
% ===========================================================================

A model \keyterm{overfits} when training loss keeps falling while validation
loss rises: it is memorising rather than
generalising\textsuperscript{\cite[][pp.~27--29]{fleuretLittleBookDeep}}.

\begin{figure}[htbp]
\centering
\begin{tikzpicture}
\begin{axis}[almanach axis,
  height=46mm,
  domain=1:60,
  samples=120,
  xlabel={epoch},
  ylabel={loss},
  ymin=0,ymax=1.5,
  legend pos=north east]
  \addplot[AlmanachTeal,thick]  {0.08+1.2*exp(-x/9)};
    \addlegendentry{training loss}
  \addplot[AlmanachOchre,thick] {0.34+1.0*exp(-x/7)+0.00055*(x-18)^2};
    \addlegendentry{validation loss}
  \draw[AlmanachRed,dashed,thick] (axis cs:22,0) -- (axis cs:22,1.5);
  \node[font=\sffamily\scriptsize,color=AlmanachRed,anchor=west]
    at (axis cs:23.5,1.28) {early-stopping point};
\end{axis}
\end{tikzpicture}
\caption{The canonical overfitting signature. Training loss decreases
monotonically while validation loss turns upward; the minimum of the validation
curve, not the end of training, identifies the checkpoint worth keeping. Note
that the training curve alone contains no evidence of the problem.}
\label{fig:dl-overfitting}
\end{figure}

\begin{table}[htbp]
\centering
\caption{Regularisers, the prior each one encodes, and what it cannot fix. No
regulariser compensates for unrepresentative data.}
\label{tab:dl-regularisers}
\begin{tblr}{
  width=\textwidth,
  colspec={Q[l,0.18\textwidth] X[l] X[l]},
  row{1}={font=\bfseries,bg=AlmanachMist},
  hlines,vlines}
Technique & Assumption it encodes & Blind spot \\
Early stopping & The best generaliser appears before convergence &
  Needs a validation set that genuinely resembles deployment \\
L2 / weight decay & Small, distributed weights generalise better &
  Does not prevent memorising through many small weights \\
L1 & Few features actually matter &
  Unstable selection among correlated features \\
Dropout\textsuperscript{\cite{srivastavaDropoutSimpleWay2014}} &
  No unit may be indispensable &
  Interacts badly with batch norm; rarely used in modern Transformers \\
Data augmentation & Known transforms preserve the label &
  A wrong invariance teaches a wrong fact \\
More data & Coverage beats constraint &
  Only if the new data covers the failing region \\
\end{tblr}
\end{table}

\subsection{Dropout, Stated Correctly}

During training, \keyterm{dropout} zeroes each activation independently with
probability $p$\textsuperscript{\cite{srivastavaDropoutSimpleWay2014}}. Units
cannot co-adapt, because no unit can rely on any particular partner being
present.

\begin{cautionbox}[Inverted dropout: what your framework actually does]
The original paper describes keeping activations unscaled during training and
multiplying by $(1-p)$ at test time. Every mainstream framework does the
opposite, called \keyterm{inverted dropout}: it divides the surviving
activations by $(1-p)$ \emph{during training} and does nothing at all at test
time. The expected value is the same, but the inference path stays a plain
forward pass, which matters for export and deployment. If you read the $(1-p)$
test-time rule in a textbook and then look for it in framework source code, you
will not find it---and the discrepancy is a common source of confusion when
re-implementing dropout by hand.
\end{cautionbox}

\begin{workedexample}[Dropout preserves the expected activation]
Let a unit output $a = 4.0$ with drop probability $p = 0.25$. Under inverted
dropout the training-time output is
\[
  \tilde a =
  \begin{cases}
    a/(1-p) = 4.0/0.75 = 5.333 & \text{with probability } 0.75,\\
    0 & \text{with probability } 0.25.
  \end{cases}
\]
Then $\mathbb{E}[\tilde a] = 0.75\times5.333 + 0.25\times0 = 4.0 = a$.
\checkmark\ The next layer therefore sees the same expected input scale whether
dropout is on or off, which is exactly why no test-time correction is needed.
\end{workedexample}

\begin{cautionbox}[Overfitting is not the only generalisation story]
Two results should stop you from over-applying the classical picture. First,
deep networks can fit \emph{random labels} perfectly, so classical capacity
measures do not explain their
generalisation\textsuperscript{\cite{zhangUnderstandingDeepLearning2021}}.
Second, test error can follow a \keyterm{double descent} curve: it rises as the
model approaches the interpolation threshold and then falls \emph{again} as the
model grows past
it\textsuperscript{\cite{nakkiranDeepDoubleDescent2020}}. The textbook
U-shaped bias--variance picture is a good first model and a poor last word.
Measure the curve for your own setting rather than assuming its shape.
\end{cautionbox}

\begin{checkpoint}
Your validation accuracy is $94\,\%$ and your test accuracy is $71\,\%$. List
three causes that regularisation would \emph{not} fix, and name the diagnostic
that distinguishes them.
\end{checkpoint}

% ===========================================================================
\section{Convolutional Networks}
% ===========================================================================

A \keyterm{convolutional neural network} replaces the dense matrix product with
a small learned \keyterm{filter} slid across the
input\textsuperscript{\cite[][pp.~43--48]{fleuretLittleBookDeep}}. Three
properties follow, and all three are consequences of the sliding, not extra
assumptions:

\begin{description}
  \item[Weight sharing] The same filter is applied at every position, so a
    feature detected in one corner is detected in all of them
    (\emph{translation equivariance}).
  \item[Local connectivity] Each output depends on a small neighbourhood, which
    encodes the prior that nearby pixels are related.
  \item[Parameter efficiency] Parameters scale with the filter, not with the
    input size.
\end{description}

\begin{plainlanguage}
A convolution is a rubber stamp with a scoring pattern on it. You press the same
stamp at every position of the image and write down how strongly the pattern
matched there. Learning a filter means learning what the stamp should look like;
using many filters means carrying a whole set of stamps for edges, textures,
and later, eyes and wheels.
\end{plainlanguage}

\begin{workedexample}[A convolution computed by hand]
Take a $3\times3$ input containing a bright vertical edge, and a $2\times2$
vertical-edge filter:
\[
  X = \begin{pmatrix}
    1 & 1 & 8\\ 1 & 1 & 8\\ 1 & 1 & 8
  \end{pmatrix},
  \qquad
  K = \begin{pmatrix} 1 & -1\\ 1 & -1 \end{pmatrix}.
\]
With stride $1$ and no padding the output is $2\times2$. Each entry is the sum
of the element-wise product over the covered window:
\[
\begin{aligned}
  Y_{11} &= 1(1) + 1(-1) + 1(1) + 1(-1) = 0, &
  Y_{12} &= 1(1) + 8(-1) + 1(1) + 8(-1) = -14,\\
  Y_{21} &= 1(1) + 1(-1) + 1(1) + 1(-1) = 0, &
  Y_{22} &= 1(1) + 8(-1) + 1(1) + 8(-1) = -14.
\end{aligned}
\]
\[
  Y = \begin{pmatrix} 0 & -14\\ 0 & -14 \end{pmatrix}.
\]
\emph{Read the result.} The filter output is zero in the flat region and
strongly negative exactly where the intensity jumps from $1$ to $8$. The filter
did not detect ``brightness''; it detected \emph{change} in the horizontal
direction, and it did so with four parameters applied identically everywhere.
The sign tells you the direction of the edge---a filter with $K$ negated would
return $+14$ on the same image.
\end{workedexample}

\begin{figure}[htbp]
\centering
\begin{tikzpicture}[
  cell/.style={draw=AlmanachInk!35,minimum size=7.5mm,font=\scriptsize,
    inner sep=0pt},
  incell/.style={cell,fill=AlmanachMist},
  kcell/.style={cell,fill=AlmanachOchre!25,draw=AlmanachOchre},
  ocell/.style={cell,fill=AlmanachTeal!18,draw=AlmanachTeal},
  lbl/.style={font=\sffamily\scriptsize,color=AlmanachInk!70},
  arr/.style={-{Stealth[length=1.8mm]},thick,draw=AlmanachOchre}]

  \foreach \r [count=\i from 0] in {{1,1,8},{1,1,8},{1,1,8}}{
    \foreach \v [count=\j from 0] in \r {
      \node[incell] at (\j*7.5mm, -\i*7.5mm) {\v};}}
  \draw[AlmanachOchre,line width=1pt]
    (-3.75mm,3.75mm) rectangle (11.25mm,-11.25mm);
  \node[lbl] at (7.5mm,-26mm) {input $X$ ($3\times3$)};
  \node[lbl,color=AlmanachOchre] at (3.75mm,9mm) {window};

  \node[lbl] at (34mm,-7.5mm) {$\ast$};

  \foreach \r [count=\i from 0] in {{1,-1},{1,-1}}{
    \foreach \v [count=\j from 0] in \r {
      \node[kcell] at (46mm+\j*7.5mm, -\i*7.5mm) {\v};}}
  \node[lbl] at (49.75mm,-26mm) {filter $K$};

  \draw[arr] (60mm,-7.5mm) -- (72mm,-7.5mm);

  \foreach \r [count=\i from 0] in {{0,-14},{0,-14}}{
    \foreach \v [count=\j from 0] in \r {
      \node[ocell,minimum size=9mm] at (80mm+\j*9mm, -\i*9mm) {\v};}}
  \node[lbl] at (84.5mm,-26mm) {output $Y$ ($2\times2$)};
\end{tikzpicture}
\caption{The worked convolution in full. The highlighted window slides across
every valid position; at each stop the filter values multiply the covered
inputs and the products are summed. The output is zero over the flat region and
$-14$ exactly at the edge, so four shared parameters have produced a
position-independent edge detector.}
\label{fig:dl-convolution}
\end{figure}

\subsection{Output Size and Parameter Count}

For a square input of side $W$, filter side $F$, padding $P$, and stride $S$:
\begin{equation}
  W_{\text{out}} = \left\lfloor \frac{W - F + 2P}{S} \right\rfloor + 1 .
  \label{eq:dl-conv-size}
\end{equation}
\where{$W$ is the input side length, $F$ the filter side length, $P$ the
padding added to each edge, and $S$ the stride. The floor function accounts for
a final window that does not fit.}

\begin{workedexample}[Shapes and the parameter saving, with numbers]
\emph{Output size.} A $32\times32$ input, $3\times3$ filter, padding $1$,
stride $1$:
\[
  W_{\text{out}} = \left\lfloor \frac{32-3+2}{1} \right\rfloor + 1
  = 31 + 1 = 32 .
\]
Padding of $1$ with a $3\times3$ filter preserves the spatial size---this is why
``same'' padding is the default in most architectures.

\emph{Parameter count.} A convolutional layer with $C_{\text{in}}$ input
channels, $C_{\text{out}}$ filters of side $F$ has
$F^2 C_{\text{in}} C_{\text{out}} + C_{\text{out}}$ parameters. For
$C_{\text{in}}=1$, $C_{\text{out}}=32$, $F=3$:
\[
  9 \times 1 \times 32 + 32 = 288 + 32 = 320 \text{ parameters.}
\]
Compare with the dense layer from the earlier example, $100{,}480$ parameters
for a $784\to128$ map. The convolution uses
\[
  100{,}480 / 320 \approx 314\times \text{ fewer parameters}
\]
while producing $32$ feature maps over the \emph{whole} image rather than a
single fixed-size vector. That ratio is the entire economic argument for
convolution on grid-structured data.
\end{workedexample}

\subsection{Pooling, Receptive Fields, and Landmark Architectures}

\keyterm{Pooling} (max or average) reduces spatial resolution and buys a degree
of translation \emph{invariance}\textsuperscript{\cite[][pp.~49--50]{fleuretLittleBookDeep}}.
Stacking convolution and pooling builds a \keyterm{feature hierarchy}: edges and
textures early, parts and objects late. The \keyterm{receptive field}---the
input region influencing one output unit---grows with depth, which is why deep
CNNs can reason about large structures using only small filters.

\begin{table}[htbp]
\centering
\caption{Landmark convolutional architectures, each cited to its original
paper. The column that matters is the third: every entry solved a specific
obstacle rather than simply being larger.}
\label{tab:dl-cnn-landmarks}
\begin{tblr}{
  width=\textwidth,
  colspec={Q[l,0.14\textwidth] Q[c,0.09\textwidth] X[l]},
  row{1}={font=\bfseries,bg=AlmanachMist},
  hlines,vlines}
Architecture & Year & Obstacle removed \\
LeNet-5\textsuperscript{\cite{lecunGradientBasedLearningApplied1998}} & 1998 &
  Showed end-to-end gradient training on images was practical at all \\
AlexNet\textsuperscript{\cite{krizhevskyImageNetClassificationDeep2012}} &
  2012 & ReLU, dropout, and GPU training made large-scale vision feasible \\
VGG\textsuperscript{\cite{simonyanVeryDeepConvolutional2015}} & 2014 &
  Demonstrated that stacked $3\times3$ filters beat larger ones at equal cost \\
ResNet\textsuperscript{\cite{heDeepResidualLearning2016}} & 2015 &
  Residual shortcuts removed the depth barrier (\cref{eq:dl-residual}) \\
\end{tblr}
\end{table}

\begin{engineernote}
Convolution encodes a prior: \emph{nearby inputs are related and the same
pattern can occur anywhere}. That prior is true for images, spectrograms, and
many time series. It is false for tabular data, where column order is
arbitrary. Applying a CNN to a spreadsheet is not a conservative default---it is
an assertion about your columns that is almost certainly wrong.
\end{engineernote}

\begin{checkpoint}
A network stacks three $3\times3$ convolutions with stride $1$ and no pooling.
What is the receptive field of one output unit, in input pixels? Compare the
parameter count with a single $7\times7$ convolution achieving the same
receptive field, and explain VGG's design choice from your two numbers.
\end{checkpoint}

% ===========================================================================
\section{Sequence Models: RNN, LSTM, and GRU}
% ===========================================================================

A \keyterm{recurrent neural network} carries a hidden state $h_t$ forward
through time\textsuperscript{\cite[][pp.~55--58]{fleuretLittleBookDeep}}:
\begin{equation}
  h_t = \phi\!\left(W_h h_{t-1} + W_x x_t + b\right).
  \label{eq:dl-rnn}
\end{equation}
\where{$x_t$ is the input at step $t$, $h_t$ the hidden state carried
forward, $W_x$ and $W_h$ the input and recurrent weight matrices, $b$ the bias,
and $\phi$ the activation. Note that $W_h$ is reused at every step --- which is
what causes both the parameter efficiency and the gradient problem.}
The same parameters handle sequences of any length, which is the appeal. The
repeated multiplication by $W_h$ is also the problem.

\begin{figure}[htbp]
\centering
\begin{tikzpicture}[
  node distance=11mm,
  cellb/.style={draw=AlmanachTeal,fill=AlmanachMist,rounded corners=1mm,
    minimum height=8mm,minimum width=12mm,font=\sffamily\scriptsize},
  io/.style={font=\scriptsize,color=AlmanachInk!75},
  arr/.style={-{Stealth[length=1.7mm]},thick,draw=AlmanachInk!55},
  rec/.style={-{Stealth[length=1.7mm]},thick,draw=AlmanachOchre}]
  \foreach \t [count=\i from 0] in {1,2,3,4}{
    \node[cellb] (c\t) at (\i*22mm,0) {$\phi$};
    \node[io,below=6mm of c\t] (x\t) {$x_\t$};
    \node[io,above=6mm of c\t] (y\t) {$h_\t$};
    \draw[arr] (x\t) -- (c\t);
    \draw[arr] (c\t) -- (y\t);}
  \foreach \a/\b in {c1/c2,c2/c3,c3/c4}{
    \draw[rec] (\a) -- (\b) node[midway,above,io,color=AlmanachOchre]
      {$W_h$};}
  \node[io,left=8mm of c1] (h0) {$h_0$};
  \draw[rec] (h0) -- (c1);
  \draw[decorate,decoration={brace,amplitude=3.5pt,mirror},AlmanachRed]
    ($(c1.south west)+(-2mm,-13mm)$) -- ($(c4.south east)+(2mm,-13mm)$)
    node[midway,below=4pt,io,color=AlmanachRed,align=center]
    {gradient is multiplied by $W_h$ once per step};
\end{tikzpicture}
\caption{An RNN unrolled over four time steps. The same weights are reused at
every step, so the network handles any sequence length---but the backward pass
multiplies by the same $W_h$ once per step, which is the exact mechanism behind
vanishing and exploding gradients.}
\label{fig:dl-rnn-unrolled}
\end{figure}

\begin{workedexample}[Vanishing and exploding gradients, quantified]
Gradients flowing back $T$ steps are scaled by roughly $\lVert W_h\rVert^T$.
Take a sequence of $T = 100$ steps:
\[
  0.9^{100} = 2.66\times10^{-5},
  \qquad
  0.99^{100} = 0.366,
  \qquad
  1.1^{100} = 1.38\times10^{4}.
\]
A spectral norm of $0.9$---only slightly contractive---already attenuates the
learning signal by a factor of $37{,}600$ across a $100$-token sentence. A norm
of $1.1$ explodes it by a comparable factor in the other direction. The window
in which plain recurrence trains stably is uncomfortably narrow, and it narrows
further as sequences lengthen. Gradient clipping treats the explosion; it does
nothing for the vanishing case.
\end{workedexample}

\subsection{LSTM: An Additive Path Through Time}

The \keyterm{LSTM} adds a \keyterm{cell state} $c_t$ and three multiplicative
gates\textsuperscript{\cite{hochreiterLongShortTermMemory1997},
\cite[][pp.~59--61]{fleuretLittleBookDeep}}:
\begin{equation}
  c_t = f_t\odot c_{t-1} + i_t\odot\tilde c_t,
  \qquad
  h_t = o_t\odot\tanh(c_t).
  \label{eq:dl-lstm}
\end{equation}
\where{$c_t$ is the cell state, $h_t$ the hidden state, $\tilde c_t$ the
candidate update, and $f_t,i_t,o_t\in(0,1)$ the forget, input, and output
gates; $\odot$ is element-wise multiplication.}

The decisive detail is the $+$ in \cref{eq:dl-lstm}. The cell state is
\emph{added to}, not repeatedly multiplied by a weight matrix, so when the
forget gate $f_t$ stays near $1$ the gradient path along $c$ is close to the
identity---the same trick as the residual connection in
\cref{eq:dl-residual}, discovered nearly two decades earlier.

\begin{figure}[htbp]
\centering
\begin{tikzpicture}[
  gate/.style={draw=AlmanachOchre,fill=AlmanachOchre!20,rounded corners=1mm,
    minimum height=6mm,minimum width=9mm,font=\sffamily\scriptsize},
  opn/.style={circle,draw=AlmanachTeal,fill=white,inner sep=0.7pt,
    font=\scriptsize},
  lbl/.style={font=\sffamily\scriptsize,color=AlmanachInk!75},
  hw/.style={-{Stealth[length=1.8mm]},line width=1.1pt,draw=AlmanachTeal},
  arr/.style={-{Stealth[length=1.6mm]},draw=AlmanachInk!55}]

  \node[lbl] (cin) at (-1.5,1.6) {$c_{t-1}$};
  \node[opn] (mul) at (0.6,1.6) {$\times$};
  \node[opn] (add) at (3.0,1.6) {$+$};
  \node[lbl] (cout) at (5.2,1.6) {$c_t$};
  \draw[hw] (cin) -- (mul);
  \draw[hw] (mul) -- (add);
  \draw[hw] (add) -- (cout);
  \node[lbl,above=1mm of add,color=AlmanachTeal,align=center]
    {additive path:\\gradient survives};

  \node[gate] (f) at (0.6,0.2) {$f_t$};
  \node[gate] (i) at (2.3,0.2) {$i_t$};
  \node[gate] (g) at (3.7,0.2) {$\tilde c_t$};
  \node[gate] (o) at (5.2,0.2) {$o_t$};
  \node[opn]  (mi) at (3.0,0.9) {$\times$};

  \draw[arr] (f) -- (mul);
  \draw[arr] (i) -- (mi); \draw[arr] (g) -- (mi);
  \draw[arr] (mi) -- (add);

  \node[opn] (th) at (5.2,0.95) {$\tanh$};
  \node[opn] (mo) at (5.2,-0.75) {$\times$};
  \draw[arr] (add.south) .. controls +(0,-0.25) and +(0,0.3) .. (th);
  \draw[arr] (th) -- (o);
  \draw[arr] (o) -- (mo);
  \node[lbl] at (6.3,-0.75) {$h_t$};
  \draw[arr] (mo) -- (6.05,-0.75);

  \node[lbl] at (0.6,-0.75) {forget};
  \node[lbl] at (3.0,-0.75) {write};
  \node[lbl,anchor=north] at (5.2,-1.15) {read};
  \node[lbl,align=center] at (1.7,-1.6)
    {inputs $x_t$ and $h_{t-1}$ feed every gate};
\end{tikzpicture}
\caption{The LSTM cell viewed as memory hardware: the forget gate decides what
to erase, the input gate what to write, and the output gate what to read. The
thick horizontal line is the cell state---an additive highway that lets
gradients travel many time steps without repeated matrix multiplication.}
\label{fig:dl-lstm-cell}
\end{figure}

The \keyterm{GRU} merges the cell and hidden state and uses two gates---a reset
gate and an update gate---achieving similar quality with fewer
parameters\textsuperscript{\cite{choLearningPhraseRepresentations2014},
\cite[][p.~62]{fleuretLittleBookDeep}}.

\begin{engineernote}
Transformers have displaced recurrence for most sequence tasks, but recurrent
models are not obsolete. They keep a genuine structural advantage where the
input arrives one element at a time and memory is bounded: an RNN carries
$O(1)$ state per step, whereas an attention model over a growing context carries
$O(T)$. For streaming keyword spotting on a microcontroller, that difference
decides the architecture.
\end{engineernote}

\begin{checkpoint}
An LSTM's forget gate saturates at $f_t \approx 1$ for a long span. Describe
what happens to (a) the information stored in $c_t$ and (b) the gradient flowing
back through that span. Then name the failure this creates if $f_t$ never
drops.
\end{checkpoint}

% ===========================================================================
\section{Attention and the Transformer}
\label{sec:dl-attention}
% ===========================================================================

Dense layers cannot handle variable-length input; convolutions move information
only locally; recurrence moves it sequentially. \keyterm{Attention} computes,
for every output position, a weighted average over the \emph{entire} input,
regardless of distance\textsuperscript{\cite[][p.~73]{fleuretLittleBookDeep}}.

\begin{plainlanguage}
Attention is a soft library lookup. Each position writes a \emph{query}---``I am
a verb, I need my subject.'' Every position also advertises a \emph{key}---``I am
a plural noun''---and holds \emph{value} content. Match each query against all
keys to get a relevance score, turn the scores into percentages, and mix the
values in those proportions. Unlike a real library, you do not get one book: you
get a blend of all of them, weighted by relevance.
\end{plainlanguage}

\subsection{Scaled Dot-Product Attention}

Given queries $Q\in\mathbb{R}^{N_q\times d_k}$, keys
$K\in\mathbb{R}^{N_{kv}\times d_k}$, and values
$V\in\mathbb{R}^{N_{kv}\times d_v}$\textsuperscript{\cite[][pp.~74--76]{fleuretLittleBookDeep}}:
\begin{equation}
  \mathrm{Attention}(Q,K,V)
  = \mathrm{softmax}\!\left(\frac{QK^{\!\top}}{\sqrt{d_k}}\right)V .
  \label{eq:dl-attention}
\end{equation}
\where{$Q$ holds one query vector per output position, $K$ one key per input
position, and $V$ one value per input position; $d_k$ is the dimension of the
query and key vectors, and dividing by $\sqrt{d_k}$ keeps the scores in a range
where softmax still has usable gradients.}
Read it right to left as four steps: score every query against every key
($QK^{\!\top}$), rescale ($/\sqrt{d_k}$), normalise into weights (softmax), then
mix the values ($\cdot V$).

\begin{workedexample}[Attention on three tokens, computed by hand]
Work in $d_k = d_v = 2$ with three tokens. Let
\[
  K = \begin{pmatrix} 1&0\\ 0&1\\ 1&1 \end{pmatrix},
  \qquad
  V = \begin{pmatrix} 2&0\\ 0&2\\ 1&1 \end{pmatrix},
\]
and take the query of the third token, $q_3 = (1,\,1)$.

\emph{Step 1---score.} Dot products against each key row:
\[
  q_3\!\cdot\!k_1 = 1,\qquad q_3\!\cdot\!k_2 = 1,\qquad q_3\!\cdot\!k_3 = 2 .
\]

\emph{Step 2---scale} by $1/\sqrt{d_k} = 1/\sqrt{2} = 0.7071$:
\[
  s = (0.7071,\ 0.7071,\ 1.4142).
\]

\emph{Step 3---softmax.} $e^{0.7071}=2.028$, $e^{1.4142}=4.113$; the sum is
$2.028+2.028+4.113 = 8.169$, so
\[
  \alpha = (0.248,\ 0.248,\ 0.503).
\]
Check: $0.248+0.248+0.503 = 0.999 \approx 1$. \checkmark

\emph{Step 4---mix the values.}
\[
  y_3 = 0.248\begin{pmatrix}2\\0\end{pmatrix}
      + 0.248\begin{pmatrix}0\\2\end{pmatrix}
      + 0.503\begin{pmatrix}1\\1\end{pmatrix}
      = \begin{pmatrix}0.497+0.503\\ 0.497+0.503\end{pmatrix}
      = \begin{pmatrix}1.00\\ 1.00\end{pmatrix}.
\]

\emph{Interpret.} Token three attends slightly more than half its budget to
itself (its key matches its own query best) and splits the rest evenly between
the other two. The output is a blend, not a selection. Every number in
\cref{eq:dl-attention} has now been exercised on real values---if you can
reproduce this table from memory, you understand attention mechanically.
\end{workedexample}

\begin{workedexample}[Why divide by $\sqrt{d_k}$?]
Repeat step 3 of the previous example \emph{without} the scaling. The raw scores
are $(1,1,2)$, giving $e^1=2.718$, $e^2=7.389$, sum $12.826$, and
\[
  \alpha_{\text{unscaled}} = (0.212,\ 0.212,\ 0.576)
  \quad\text{versus}\quad
  \alpha_{\text{scaled}} = (0.248,\ 0.248,\ 0.503).
\]
The unscaled distribution is noticeably peakier. Now consider why this gets
worse with dimension: if query and key components are independent with unit
variance, the dot product $q\!\cdot\!k$ has variance $d_k$, so its typical
magnitude grows like $\sqrt{d_k}$. At $d_k = 64$ the raw scores would spread
about eight times wider than at $d_k = 1$, driving softmax into a near-one-hot
regime where its gradient is almost zero. Dividing by $\sqrt{d_k}$ cancels that
growth exactly and keeps the softmax in a trainable
range\textsuperscript{\cite[][p.~4]{vaswaniAttentionAllYou2023}}.
\end{workedexample}

\begin{figure}[htbp]
\centering
\begin{tikzpicture}[
  cell/.style={draw=AlmanachInk!30,minimum size=9mm,font=\scriptsize,
    inner sep=0pt},
  lbl/.style={font=\sffamily\scriptsize,color=AlmanachInk!75},
  hd/.style={font=\sffamily\scriptsize,color=AlmanachTeal}]

  \node[hd] at (13.5mm,12mm) {attention weights $\alpha$};
  \foreach \r [count=\i from 0] in {
    {1.00/100,0.00/0,0.00/0},
    {0.27/27,0.73/73,0.00/0},
    {0.25/25,0.25/25,0.50/50}}{
    \foreach \pair [count=\j from 0] in \r {
      \foreach \v/\s in \pair {
        \node[cell,fill=AlmanachTeal!\s] at (\j*9mm,-\i*9mm) {\v};}}}
  \foreach \t [count=\i from 0] in {the,cat,sat}{
    \node[lbl,left=1mm] at (-4.5mm,-\i*9mm) {\t};}
  \foreach \t [count=\j from 0] in {the,cat,sat}{
    \node[lbl,above=1mm] at (\j*9mm,4.5mm) {\t};}
  \node[lbl,rotate=90] at (-15mm,-9mm) {query};
  \node[lbl] at (13.5mm,-33mm) {key};

  \node[lbl,align=left,anchor=west,text width=52mm] at (36mm,-9mm)
    {The upper triangle is \emph{masked} to zero: when predicting a token,
    the model may not look at tokens that come after it. Each row sums to
    one, so every query spends exactly one unit of attention budget.};
\end{tikzpicture}
\caption{A causal attention matrix for a three-token sequence. Row $i$ shows how
query $i$ distributes its attention over the keys; the enforced zeros above the
diagonal are what make a decoder autoregressive rather than clairvoyant.}
\label{fig:dl-attention-matrix}
\end{figure}

\subsection{Self-Attention, Cross-Attention, and Multiple Heads}

In \keyterm{self-attention} the queries, keys, and values are all projections of
the same sequence, so every position can gather from every
other\textsuperscript{\cite[][p.~79]{fleuretLittleBookDeep}}. In
\keyterm{cross-attention} the keys and values come from a different sequence,
which is how a decoder reads an encoder.

A single head averages over all positions and cannot track several kinds of
relationship at once. \keyterm{Multi-head attention} runs $H$ heads in
parallel, each with its own projections, then concatenates and
projects\textsuperscript{\cite[][pp.~4--5]{vaswaniAttentionAllYou2023}}:
\begin{equation}
  \mathrm{MultiHead}(Q,K,V)
  = \mathrm{Concat}(\mathrm{head}_1,\dots,\mathrm{head}_H)\,W^O .
  \label{eq:dl-multihead}
\end{equation}
\where{$H$ is the number of heads, each computing
\cref{eq:dl-attention} on its own projection of the input, and $W^O$ the output
projection that mixes the concatenated heads back to the model dimension.}

\begin{workedexample}[Multi-head bookkeeping]
Take $d_{\text{model}} = 512$ and $H = 8$ heads. Each head works in
$d_k = d_v = 512/8 = 64$ dimensions. The projections cost
\[
  \underbrace{3 \times 512 \times 512}_{W^Q, W^K, W^V}
  + \underbrace{512 \times 512}_{W^O}
  = 4 \times 262{,}144 = 1{,}048{,}576 \text{ parameters.}
\]
Splitting into heads therefore costs \emph{nothing} in parameters or
computation: the same $512$ dimensions are partitioned, not duplicated. What is
bought is representational---eight separate score matrices instead of one
averaged view. This is why the head count is a cheap knob and the model
dimension is an expensive one.
\end{workedexample}

\begin{figure}[htbp]
\centering
\begin{tikzpicture}[
  node distance=4.5mm,
  blk/.style={draw=AlmanachTeal,fill=AlmanachMist,rounded corners=1mm,
    minimum height=6.5mm,minimum width=17mm,font=\sffamily\scriptsize,
    align=center},
  hb/.style={blk,fill=AlmanachOchre!22,draw=AlmanachOchre,minimum width=15mm},
  arr/.style={-{Stealth[length=1.6mm]},draw=AlmanachInk!55},
  lbl/.style={font=\sffamily\scriptsize,color=AlmanachInk!70}]
  \node[blk] (in) {input $X$\\\scriptsize $(T,512)$};
  \node[hb,above right=6mm and 6mm of in] (h2) {head 2\\\scriptsize $(T,64)$};
  \node[hb,above=of h2] (h1) {head 1\\\scriptsize $(T,64)$};
  \node[hb,below=of h2] (hd) {$\dots$};
  \node[hb,below=of hd] (h8) {head 8\\\scriptsize $(T,64)$};
  \node[blk,right=24mm of h2,yshift=-6mm] (cc)
    {concat\\\scriptsize $(T,512)$};
  \node[blk,right=of cc] (wo) {$W^O$\\\scriptsize $(T,512)$};
  \foreach \h in {h1,h2,hd,h8}{\draw[arr] (in) -- (\h); \draw[arr] (\h) -- (cc);}
  \draw[arr] (cc) -- (wo);
  \node[lbl,below=8mm of in,align=center,text width=70mm]
    {$8\times64 = 512$: the heads partition the model dimension,
     they do not enlarge it};
\end{tikzpicture}
\caption{Multi-head attention splits the model dimension across heads and
reassembles it. Because $H \times d_k = d_{\text{model}}$, adding heads changes
what the layer can represent without changing what it costs.}
\label{fig:dl-multihead}
\end{figure}

\subsection{Positional Encoding}

Attention is permutation-invariant: shuffle the tokens and
\cref{eq:dl-attention} returns the same multiset of outputs. Order must be
injected explicitly\textsuperscript{\cite[][pp.~81--82]{fleuretLittleBookDeep}}.
The original Transformer adds a sinusoidal
encoding\textsuperscript{\cite[][p.~6]{vaswaniAttentionAllYou2023}}:
\begin{equation}
  \mathrm{PE}_{t,2i}   = \sin\!\left(\frac{t}{10000^{2i/d}}\right),
  \qquad
  \mathrm{PE}_{t,2i+1} = \cos\!\left(\frac{t}{10000^{2i/d}}\right).
  \label{eq:dl-posenc}
\end{equation}
\where{$t$ is the token's position in the sequence, $d$ the model dimension,
and $i$ indexes dimension pairs; each pair oscillates at its own frequency, so
low $i$ gives fast channels and high $i$ slow ones.}

\begin{figure}[htbp]
\centering
\begin{tikzpicture}
\begin{axis}[almanach axis,
  height=42mm,
  domain=0:60,
  samples=241,
  xlabel={token position $t$},
  ylabel={$\mathrm{PE}_{t,2i}$},
  ymin=-1.3,ymax=1.7,
  legend pos=north east,
  legend columns=3]
  \addplot[AlmanachTeal,thick]  {sin(deg(x/1))};    \addlegendentry{$i=0$}
  \addplot[AlmanachOchre,thick] {sin(deg(x/6.3))};  \addlegendentry{$i=2$}
  \addplot[AlmanachRed,thick]   {sin(deg(x/40))};   \addlegendentry{$i=4$}
\end{axis}
\end{tikzpicture}
\caption{Sinusoidal position channels at three frequencies. Fast channels
distinguish neighbouring tokens; slow channels encode coarse position in the
sequence. Together they give every position a unique signature whose
\emph{differences} vary smoothly with distance, which is the property attention
can exploit.}
\label{fig:dl-positional-encoding}
\end{figure}

\begin{cautionbox}[Length extrapolation is empirical, not guaranteed]
Because \cref{eq:dl-posenc} is a deterministic formula, it \emph{can} be
evaluated at positions never seen in training. That is often misreported as a
guarantee that the model will work on longer sequences. It is not. The encoding
is defined there; whether the trained attention layers behave sensibly there is
an empirical question, and in practice quality degrades beyond the trained
context length. Learned positional embeddings performed comparably in the
original experiments\textsuperscript{\cite[][p.~6]{vaswaniAttentionAllYou2023}},
and modern models use rotary or relative schemes precisely because this problem
is unsolved by sinusoids alone. Always re-evaluate at the length you intend to
deploy.
\end{cautionbox}

\subsection{The Transformer Block}

\begin{figure}[htbp]
\centering
\begin{tikzpicture}[
  node distance=0.25cm and 0.7cm,
  box/.style={draw=AlmanachTeal,rounded corners=2pt,minimum width=2.8cm,
    minimum height=0.55cm,font=\sffamily\scriptsize,align=center,
    fill=AlmanachMist},
  addnorm/.style={draw=AlmanachOchre,rounded corners=2pt,minimum width=2.8cm,
    minimum height=0.45cm,font=\sffamily\scriptsize,align=center,
    fill=AlmanachOchre!18},
  lbl/.style={font=\sffamily\scriptsize,color=AlmanachInk!65},
  arr/.style={-{Stealth[length=2mm]},thick,draw=AlmanachInk!55}]
  \node[box] (emb_e) {Embedding};
  \node[box,above=of emb_e] (pe_e) {$+$ Pos.\ Encoding};
  \node[addnorm,above=of pe_e] (an1_e) {Add \& Norm};
  \node[box,above=of an1_e] (mha_e) {Multi-Head\\Self-Attention};
  \node[addnorm,above=of mha_e] (an2_e) {Add \& Norm};
  \node[box,above=of an2_e] (ffn_e) {Feed Forward};
  \node[addnorm,above=of ffn_e] (an3_e) {Add \& Norm};
  \node[lbl,left=0.15cm of mha_e] {$\times N$};

  \node[box,right=2.5cm of emb_e] (emb_d) {Embedding};
  \node[box,above=of emb_d] (pe_d) {$+$ Pos.\ Encoding};
  \node[addnorm,above=of pe_d] (an1_d) {Add \& Norm};
  \node[box,above=of an1_d] (cmha_d) {Masked\\Self-Attention};
  \node[addnorm,above=of cmha_d] (an2_d) {Add \& Norm};
  \node[box,above=of an2_d] (xatt_d) {Cross-Attention};
  \node[addnorm,above=of xatt_d] (an3_d) {Add \& Norm};
  \node[box,above=of an3_d] (ffn_d) {Feed Forward};
  \node[addnorm,above=of ffn_d] (an4_d) {Add \& Norm};
  \node[lbl,right=0.15cm of cmha_d] {$\times N$};
  \node[box,above=0.3cm of an4_d,fill=AlmanachTeal!22] (out)
    {Linear $+$ Softmax};

  \foreach \a/\b in {emb_e/pe_e,pe_e/an1_e,an1_e/mha_e,mha_e/an2_e,
    an2_e/ffn_e,ffn_e/an3_e}{\draw[arr] (\a) -- (\b);}
  \foreach \a/\b in {emb_d/pe_d,pe_d/an1_d,an1_d/cmha_d,cmha_d/an2_d,
    an2_d/xatt_d,xatt_d/an3_d,an3_d/ffn_d,ffn_d/an4_d,an4_d/out}{
    \draw[arr] (\a) -- (\b);}
  \draw[arr,dashed,draw=AlmanachOchre] (an3_e.east)
    .. controls +(0.4,0.8) and +(-0.4,0) .. (xatt_d.west)
    node[midway,above,font=\sffamily\tiny,color=AlmanachOchre] {$K$, $V$};
  \node[lbl,below=0.1cm of emb_e] {Source tokens};
  \node[lbl,below=0.1cm of emb_d] {Target tokens (shifted right)};
  \node[lbl,above=0.1cm of out] {Output probabilities};
\end{tikzpicture}
\caption{The original encoder--decoder Transformer. The encoder reads the whole
source in parallel; the decoder generates one token at a time, masked so it
cannot see the future, and reaches into the encoder through cross-attention.
Every sub-layer is wrapped in the residual-plus-normalisation pattern of
\cref{eq:dl-residual}.}
\label{fig:dl-transformer}
\end{figure}

\begin{table}[htbp]
\centering
\caption{Cost comparison per layer, for sequence length $T$ and width $d$. The
Transformer buys a constant path length with a quadratic term in $T$---the
trade that defines modern sequence modelling.}
\label{tab:dl-seq-complexity}
\begin{tblr}{
  width=\textwidth,
  colspec={Q[l,0.16\textwidth] Q[c,0.20\textwidth] Q[c,0.20\textwidth] X[l]},
  row{1}={font=\bfseries,bg=AlmanachMist},
  hlines,vlines}
Layer & Compute per layer & Sequential steps & Longest path between positions \\
Recurrent & $O(T d^2)$ & $O(T)$ & $O(T)$ \\
Convolutional & $O(F T d^2)$ & $O(1)$ & $O(\log_F T)$ \\
Self-attention & $O(T^2 d)$ & $O(1)$ & $O(1)$ \\
\end{tblr}
\end{table}

\subsection{Encoder-Only and Decoder-Only Descendants}

Most deployed models drop half of \cref{fig:dl-transformer}. \keyterm{BERT} is
encoder-only and bidirectional, trained with masked language
modelling\textsuperscript{\cite{devlinBERTPretrainingDeep2019}};
\keyterm{GPT} is decoder-only and causal, trained on next-token
prediction\textsuperscript{\cite{radfordImprovingLanguageUnderstanding2018}}.
The distinction is not stylistic: bidirectional models see the whole input and
suit classification and retrieval, while causal models can generate.

\begin{practicallab}[Attention without a GPU]
This lab needs a spreadsheet or twenty lines of NumPy---no model, no training.

\textbf{Goal.} Convert \cref{eq:dl-attention} from a formula you have read into
one you can predict.

\textbf{Steps.}
\begin{enumerate}
  \item Reproduce the three-token worked example numerically and confirm
        $y_3 = (1.00, 1.00)$ to two decimals.
  \item Compute $y_1$ and $y_2$ as well, then assemble the full $3\times3$
        weight matrix. Verify every row sums to $1$.
  \item Apply a causal mask: set the upper triangle of the score matrix to
        $-\infty$ \emph{before} the softmax. Confirm you reproduce the pattern
        in \cref{fig:dl-attention-matrix}, and explain why masking before
        rather than after the softmax is the only correct order.
  \item Multiply all keys and queries by $10$. Recompute $\alpha$ and record
        how peaked it becomes. You have just simulated what happens when
        $\sqrt{d_k}$ scaling is omitted.
  \item Set $V$ equal to the identity rows. What does the output tell you about
        the attention weights themselves?
\end{enumerate}

\textbf{Evidence to keep.} A table of the four weight matrices from steps 2--4
and one sentence per step on what changed and why.

\textbf{Success criterion.} You can predict the direction of change in
$\alpha$---peakier or flatter---before running each step. Predicting beats
observing; an experiment you cannot predict is a demonstration, not a test.
\end{practicallab}

\begin{checkpoint}
A sequence of $2{,}048$ tokens is processed by a self-attention layer with
$d=512$. Using \cref{tab:dl-seq-complexity}, compute the ratio of attention
cost to feed-forward cost. At which sequence length do the two become equal?
\end{checkpoint}

% ===========================================================================
\section{Beyond the Dense Transformer}
\label{sec:dl-beyond-transformer}
% ===========================================================================

\Cref{tab:dl-seq-complexity} named the Transformer's price: attention costs
$O(T^2 d)$ in sequence length, and every parameter is used for every token.
Two independent lines of work attack those two costs, and current frontier
models generally use both.

\begin{figure}[htbp]
\centering
\begin{tikzpicture}[
  ax/.style={->,draw=AlmanachInk!55,>={Stealth[length=1.7mm]}},
  bx/.style={draw=AlmanachTeal,fill=AlmanachMist,rounded corners=1mm,
    minimum height=8mm,minimum width=26mm,font=\sffamily\scriptsize,
    align=center},
  lbl/.style={font=\sffamily\scriptsize,color=AlmanachInk!75}]
  \draw[ax] (0,0) -- (86mm,0)
    node[right,lbl,align=left,text width=22mm] {cheaper per\\token};
  \draw[ax] (0,0) -- (0,42mm)
    node[above,lbl,align=center,text width=26mm] {more capacity\\per FLOP};
  \node[bx] at (22mm,6mm) {dense Transformer\\\scriptsize $O(T^2)$, all
    parameters};
  \node[bx,fill=AlmanachOchre!20,draw=AlmanachOchre] at (22mm,32mm)
    {sparse MoE\\\scriptsize few experts active};
  \node[bx] at (68mm,6mm) {SSM / linear\\\scriptsize $O(T)$, fixed state};
  \node[bx,fill=AlmanachTeal!25] at (68mm,32mm)
    {hybrid $+$ MoE\\\scriptsize e.g.\ Jamba};
  \node[lbl,anchor=west,align=left,text width=40mm] at (0,-9mm)
    {the two axes are independent --- a model can move along either, or both};
\end{tikzpicture}
\caption{Two independent economies. Moving right reduces the cost of
\emph{sequence length}; moving up reduces the cost of \emph{parameter count}.
Mixture-of-experts and state-space methods solve different problems and
compose, which is why recent architectures combine them.}
\label{fig:dl-architecture-axes}
\end{figure}

\subsection{Mixture of Experts: Decoupling Capacity from Compute}

In a dense network every parameter participates in every forward pass, so
capacity and compute are locked together. A \keyterm{mixture-of-experts} (MoE)
layer breaks that link: it holds $N$ parallel expert sub-networks and a
\keyterm{router} that sends each token to only the top $k$ of
them\textsuperscript{\cite{shazeerOutrageouslyLargeNeural2017}}. Typically the
MLP block of each Transformer layer is replaced.

\begin{equation}
  y = \sum_{i \in \mathrm{TopK}(g(x))} g_i(x)\, E_i(x),
  \qquad
  g(x) = \mathrm{softmax}\bigl(W_g x\bigr),
  \label{eq:dl-moe}
\end{equation}
where $E_i$ is the $i$-th expert and $g$ the routing distribution. Because
only $k$ of $N$ experts run, compute scales with $k$ while capacity scales
with $N$.

\begin{plainlanguage}
A dense model is a single generalist who personally answers every question. A
mixture-of-experts model is a large firm with a receptionist: the firm employs
hundreds of specialists, but any one question is routed to two of them. The
payroll is enormous; the cost of answering one question is not. What you save
is time, not salary --- and that is exactly the trade, because all the experts
must still be held in memory.
\end{plainlanguage}

\begin{workedexample}[Total against active parameters]
Mixtral 8$\times$7B has eight experts per layer and routes each token to two of
them\textsuperscript{\cite{jiangMixtralExperts2024}}. Its parameter counts are
often misread, so separate them:

\smallskip
\begin{tblr}{
  width=0.86\linewidth,
  colspec={Q[l,0.30\linewidth] Q[c] X[l]},
  row{1}={font=\bfseries,bg=AlmanachMist},
  hlines,vlines}
Quantity & Value & Determines \\
Total parameters & ${\approx}46.7$B & Memory to load the model \\
Active per token & ${\approx}12.9$B & Compute and latency per token \\
Experts per layer & 8 & Capacity \\
Routed per token & 2 & Cost \\
\end{tblr}
\smallskip

\emph{Read the consequence.} You must provision memory for $46.7$B parameters
but pay compute for roughly $12.9$B --- so the model has the quality profile of
something much larger than its speed suggests, and the hardware footprint of
something much larger than its speed suggests too. \emph{Both} statements are
true, and vendors tend to quote whichever flatters.

\emph{Why the name ``8$\times$7B'' misleads.} $8 \times 7 = 56 \neq 46.7$,
because only the MLP blocks are replicated across experts; attention and
embeddings are shared. Never infer memory from an MoE model's name.
\end{workedexample}

\begin{cautionbox}[Routing is a load-balancing problem, and it can collapse]
Nothing in \cref{eq:dl-moe} prevents the router from sending most tokens to a
few popular experts. Left alone it usually does --- a rich-get-richer loop in
which a slightly better expert receives more tokens, trains faster, and becomes
more attractive still. The result is \keyterm{expert collapse}: nominal
capacity of $N$ experts, effective capacity of two or three.

The standard remedy is an \keyterm{auxiliary load-balancing loss} that
penalises uneven routing, plus a per-expert capacity factor that drops or
reroutes overflow tokens\textsuperscript{\cite{fedusSwitchTransformersScaling2022}}.
Both introduce their own artefacts: token dropping makes the forward pass
depend on which \emph{other} tokens happen to share the batch, which breaks the
usual assumption that inference is per-example deterministic. If you are
evaluating an MoE model, fix the batch composition before comparing runs.
\end{cautionbox}

\subsection{State-Space Models: Linear Cost in Sequence Length}

Attention compares every token with every other token, which is what makes it
both powerful and quadratic. \keyterm{State-space models} (SSMs) return to a
recurrent formulation --- but one designed so that training can still be
parallelised.

A continuous linear system, discretised, gives
\begin{equation}
  h_t = \bar{A}h_{t-1} + \bar{B}x_t,
  \qquad
  y_t = C h_t ,
  \label{eq:dl-ssm}
\end{equation}
\where{$x_t$ is the input at step $t$, $h_t$ the hidden state, $y_t$ the
output, and $\bar A,\bar B,C$ the discretised system matrices. In S4 these are
fixed; Mamba's contribution is making $\bar B$ and $C$ depend on the input.}
which is an RNN at inference time --- $O(1)$ state per step, $O(T)$ overall ---
but, because it is \emph{linear} in $h$, it can also be evaluated as a global
convolution and therefore trained in
parallel\textsuperscript{\cite{guEfficientlyModelingLong2022}}.

\keyterm{Mamba} adds the missing ingredient:
\keyterm{selectivity}\textsuperscript{\cite{guMambaLinearTimeSequence2024}}.
In S4 the matrices are fixed, so the model compresses history the same way
regardless of content. Mamba makes $\bar B$, $C$, and the discretisation step
\emph{functions of the input}, letting the model choose what to keep and what
to forget --- recovering much of what attention provides. Because that breaks
the convolution trick, Mamba supplies a hardware-aware parallel scan instead.

\begin{table}[htbp]
\centering
\caption{The three sequence-mixing families compared on the properties that
decide deployment. The last row is the one that matters most in practice and is
quoted least often.}
\label{tab:dl-sequence-families}
\begin{tblr}{
  width=\textwidth,
  colspec={Q[l,0.21\textwidth] X[c] X[c] X[c]},
  row{1}={font=\bfseries,bg=AlmanachMist},
  hlines,vlines}
& Attention & SSM (Mamba) & Linear/RNN-style (RWKV) \\
Training cost in $T$ & $O(T^2)$ & $O(T)$ & $O(T)$ \\
Inference state & grows with $T$ (KV cache) & fixed size & fixed size \\
Per-token latency & rises with context & constant & constant \\
Exact recall from context & strong & weaker & weaker \\
\end{tblr}
\end{table}

\begin{engineernote}
The honest summary as of this book's date: SSMs win decisively on long-sequence
cost and lose on precise retrieval from context. A fixed-size state must
compress everything it has seen, so tasks that require quoting an exact string
from 50{,}000 tokens ago --- the ``needle in a haystack'' family --- favour
attention, which keeps every key and value.

That asymmetry is why the dominant production answer is neither pure
architecture but a \keyterm{hybrid}: mostly SSM layers for cheap sequence
mixing, with attention layers interleaved at intervals to preserve exact
recall, and MoE on top for
capacity\textsuperscript{\cite{lieberJambaHybridTransformerMamba2024}}. Treat
any claim that one family has simply won as a claim to check against a
benchmark that includes long-context retrieval.
\end{engineernote}

\subsection{Cheaper Attention, Without Replacing It}

Between ``keep dense attention'' and ``replace it'' sits a set of
modifications that keep the mechanism and reduce its cost. These are the
changes most likely to appear in a model you actually deploy.

\begin{table}[htbp]
\centering
\caption{Attention-preserving efficiency techniques. Note which column each one
attacks: several reduce \emph{memory} rather than compute, which is usually the
binding constraint at inference (\cref{ch:practical-llm-engineering}).}
\label{tab:dl-attention-variants}
\begin{tblr}{
  width=\textwidth,
  colspec={Q[l,0.22\textwidth] X[l] X[l]},
  row{1}={font=\bfseries,bg=AlmanachMist},
  hlines,vlines}
Technique & Mechanism & What it costs \\
Multi-query / grouped-query attention\textsuperscript{\cite{shazeerFastTransformerDecoding2019,ainslieGQATrainingGeneralized2023}}
  & Several query heads share one key/value head, shrinking the KV cache
  & A small quality loss; GQA is the compromise between MQA and full MHA \\
Sliding-window / local attention\textsuperscript{\cite{beltagyLongformerLongDocumentTransformer2020}}
  & Each token attends only within a window, plus a few global tokens
  & Long-range links must be built up across layers \\
FlashAttention\textsuperscript{\cite{daoFlashAttentionFastMemory2022}}
  & Tiles the computation to avoid writing the $T\times T$ matrix to memory
  & Nothing in accuracy --- it is exact --- but it is kernel-specific \\
Multi-head latent attention\textsuperscript{\cite{deepseekaiDeepSeekV3TechnicalReport2024}}
  & Compresses keys and values into a low-rank latent cache
  & Extra projection work, for a large cache reduction \\
Rotary position embedding\textsuperscript{\cite{suRoFormerEnhancedTransformer2024}}
  & Encodes position by rotating query and key vectors
  & Now the default; extending beyond trained length still needs care \\
Speculative decoding\textsuperscript{\cite{leviathanFastInferenceTransformers2023}}
  & A small draft model proposes tokens; the large model verifies in parallel
  & Output is \emph{identical} to the large model; gains depend on draft
  agreement \\
\end{tblr}
\end{table}

\begin{workedexample}[What grouped-query attention saves]
Take the key--value cache calculation from the deep dive later in this chapter,
but vary only the number of key/value heads. With $L=32$ layers, head dimension
$d_h=128$, \texttt{float16}, and a $8{,}192$-token context:

\[
  \text{cache} = 2 \times L \times H_{kv} \times d_h \times 2\,\text{bytes}
  \times T .
\]

\smallskip
\begin{tblr}{
  width=0.9\linewidth,
  colspec={Q[l,0.30\linewidth] Q[c] Q[c]},
  row{1}={font=\bfseries,bg=AlmanachMist},
  hlines,vlines}
Scheme & $H_{kv}$ & Cache at 8{,}192 tokens \\
Multi-head (MHA) & 32 & $4.29$ GB \\
Grouped-query (GQA) & 8 & $1.07$ GB \\
Multi-query (MQA) & 1 & $0.13$ GB \\
\end{tblr}
\smallskip

\emph{Read it.} Going from $32$ to $8$ key/value heads cuts the cache by
$4\times$ for a small, measurable quality cost --- which is why GQA is now
near-universal. Going all the way to one head saves $32\times$ but costs more
quality, so MQA appears mainly where memory is the hard constraint.

\emph{And note what did not change.} The parameter count barely moves; only the
\emph{cache} shrinks. This is a memory optimisation, and it raises the batch
size you can serve rather than the speed of any single request.
\end{workedexample}

\begin{cautionbox}[Architecture claims date fastest of anything in this book]
This section describes a landscape in motion. Specific models named here will
be superseded; the parameter counts and the relative standing of SSMs against
attention will change; and a paradigm currently absent may dominate.

What is stable is the \emph{structure} of the trade-offs, and that is what to
carry forward: capacity can be decoupled from compute, but only by paying in
memory; sequence cost can be made linear, but only by compressing history into
a fixed state, which costs exact recall; and every efficiency technique in
\cref{tab:dl-attention-variants} reduces one resource by spending another. When
you meet a new architecture, ask which of those three trades it is making
rather than whether it is better.
\end{cautionbox}

\begin{checkpoint}
A vendor offers a ``$120$B parameter'' MoE model that is ``as fast as a $15$B
dense model''. State what both numbers most likely refer to, what you must
provision, and the single question you would ask to size the hardware. Then say
which benchmark family you would insist on seeing before accepting a
long-context claim for an SSM-based model.
\end{checkpoint}

% ===========================================================================
\section{Generative Architectures}
% ===========================================================================

\subsection{Generative Adversarial Networks}

A \keyterm{GAN} trains two networks against each
other\textsuperscript{\cite{goodfellowGenerativeAdversarialNets2014},
\cite[][p.~125]{fleuretLittleBookDeep}}. The \keyterm{generator} $G(z)$ maps
noise $z\sim p(z)$ to a sample; the \keyterm{discriminator} $D(x)$ estimates the
probability that its input is real. They optimise opposite objectives:
\begin{equation}
  \min_G \max_D\ \Bigl[\,
    \mathbb{E}_{x\sim p_{\text{data}}}\bigl[\log D(x)\bigr]
    + \mathbb{E}_{z\sim p(z)}\bigl[\log\bigl(1 - D(G(z))\bigr)\bigr]
  \,\Bigr].
  \label{eq:dl-gan}
\end{equation}
\where{$G$ is the generator and $D$ the discriminator, $z$ a noise vector
drawn from the prior $p(z)$, and $x$ a real sample from the data distribution
$p_{\text{data}}$. $D(x)\in(0,1)$ is $D$'s estimated probability that its input
is real.}
At the theoretical optimum $G$ reproduces the data distribution and $D$ outputs
$0.5$ everywhere---it can do no better than guessing.

\begin{figure}[htbp]
\centering
\begin{tikzpicture}[
  node distance=8mm and 13mm,
  blk/.style={draw=AlmanachTeal,fill=AlmanachMist,rounded corners=1mm,
    minimum height=8mm,minimum width=20mm,font=\sffamily\scriptsize,
    align=center},
  gen/.style={blk,fill=AlmanachOchre!22,draw=AlmanachOchre},
  arr/.style={-{Stealth[length=1.8mm]},thick,draw=AlmanachInk!55},
  gr/.style={-{Stealth[length=1.8mm]},thick,draw=AlmanachRed,dashed},
  lbl/.style={font=\sffamily\scriptsize,color=AlmanachInk!70}]
  \node[blk] (z) {noise $z$};
  \node[gen,right=of z] (g) {Generator\\$G$};
  \node[blk,right=of g] (fake) {fake $G(z)$};
  \node[blk,below=14mm of fake] (real) {real $x$};
  \node[blk,right=of fake,yshift=-7mm] (d) {Discriminator\\$D$};
  \node[blk,right=of d] (p) {$P(\text{real})$};
  \draw[arr] (z) -- (g); \draw[arr] (g) -- (fake);
  \draw[arr] (fake) -- (d); \draw[arr] (real) -- (d); \draw[arr] (d) -- (p);
  \draw[gr] (p.north) .. controls +(0,1.4) and +(0,1.4) .. (g.north)
    node[midway,above,lbl,color=AlmanachRed]
    {gradient: ``which cue gave you away?''};
  \node[lbl,below=1mm of d,align=center,text width=32mm]
    {maximises accuracy};
  \node[lbl,below=1mm of g,align=center,text width=26mm]
    {minimises $D$'s accuracy};
\end{tikzpicture}
\caption{The adversarial loop. The discriminator's gradient tells the generator
exactly which cues betrayed the fake, and the generator removes them. Training
succeeds only while neither player is decisively winning---which is why GAN
training is delicate rather than merely slow.}
\label{fig:dl-gan}
\end{figure}

\begin{cautionbox}[Mode collapse and the balance problem]
GAN training is famously unstable, in two specific ways.
\emph{Imbalance:} if $D$ becomes too strong, $\log(1-D(G(z)))$ saturates and
the generator receives almost no gradient; if $D$ is too weak, its signal
carries no information. \emph{Mode collapse:} the generator discovers a narrow
family of outputs that reliably fools $D$ and abandons the rest of the data
distribution---the samples look excellent and the diversity is gone. Note that
neither failure shows up in the loss curves, which is what makes them
dangerous. You must evaluate coverage explicitly. Wasserstein
GANs\textsuperscript{\cite{arjovskyWassersteinGAN2017}} and spectral
normalisation address the first problem directly by supplying a gradient that
does not saturate.
\end{cautionbox}

\subsection{Autoencoders and Variational Autoencoders}

An \keyterm{autoencoder} pairs an encoder $f_\phi$ with a decoder $g_\theta$
through a narrow \keyterm{bottleneck}\textsuperscript{\cite[][p.~125]{fleuretLittleBookDeep}}.
The bottleneck is the whole design: forcing $x$ through a low-dimensional $z$
means only reconstructable structure can survive.

\begin{figure}[htbp]
\centering
\begin{tikzpicture}[
  lbl/.style={font=\sffamily\scriptsize,color=AlmanachInk!70},
  arr/.style={-{Stealth[length=1.8mm]},thick,draw=AlmanachInk!55}]
  \fill[AlmanachMist,draw=AlmanachTeal,rounded corners=1mm]
    (0,-1.1) -- (1.9,-0.42) -- (1.9,0.42) -- (0,1.1) -- cycle;
  \node[lbl] at (0.95,0) {encoder};
  \node[draw=AlmanachOchre,fill=AlmanachOchre!25,rounded corners=1mm,
    minimum width=8mm,minimum height=9mm,font=\sffamily\scriptsize]
    (z) at (2.5,0) {$z$};
  \fill[AlmanachMist,draw=AlmanachTeal,rounded corners=1mm]
    (3.1,-0.42) -- (5.0,-1.1) -- (5.0,1.1) -- (3.1,0.42) -- cycle;
  \node[lbl] at (4.05,0) {decoder};
  \node[lbl,left=2mm] at (0,0) {$x$};
  \node[lbl,right=2mm] at (5.0,0) {$\hat x$};
  \draw[arr] (-0.75,0) -- (-0.1,0);
  \draw[arr] (5.1,0) -- (5.75,0);
  \node[lbl,below=2mm of z,color=AlmanachOchre,align=center]
    {bottleneck\\$\dim z \ll \dim x$};
  \draw[arr,draw=AlmanachRed,dashed] (5.9,-0.35)
    .. controls +(0,-1.1) and +(0,-1.1) .. (-0.9,-0.35)
    node[midway,below,lbl,color=AlmanachRed]
    {reconstruction loss $\lVert x-\hat x\rVert^2$};
\end{tikzpicture}
\caption{The autoencoder. Because the only training signal is reconstruction,
the bottleneck must retain whatever is needed to rebuild the input and may
discard everything else---which is what turns it into a learned compressor,
denoiser, or anomaly detector.}
\label{fig:dl-autoencoder}
\end{figure}

A \keyterm{denoising autoencoder} receives a corrupted $\tilde x$ but targets
the clean $x$, which forces representations robust to the corruption used.

The \keyterm{variational autoencoder} makes the latent space a distribution
rather than a point\textsuperscript{\cite{kingmaAutoEncodingVariationalBayes2014}}.
The encoder emits $(\mu,\sigma)$; the \keyterm{reparameterisation trick}
$z = \mu + \sigma\odot\varepsilon$ with $\varepsilon\sim\mathcal{N}(0,I)$ moves
the randomness into an input so gradients can flow through the sampling step.
Training maximises the evidence lower bound:
\begin{equation}
  \mathcal{L}_{\text{ELBO}}
  = \mathbb{E}_{q_\phi(z\mid x)}\bigl[\log p_\theta(x\mid z)\bigr]
  - \mathrm{KL}\bigl(q_\phi(z\mid x)\,\|\,p(z)\bigr).
  \label{eq:dl-elbo}
\end{equation}
\where{$x$ is the observation, $z$ the latent code, $q_\phi(z\mid x)$ the
encoder with parameters $\phi$, and $p_\theta(x\mid z)$ the decoder with
parameters $\theta$; $p(z)$ is the prior, usually $\mathcal{N}(0,I)$. The first
term rewards reconstruction, the second keeps the latent distribution close to
the prior.}
The first term rewards reconstruction; the second pulls the latent distribution
towards $\mathcal{N}(0,I)$ so that decoding a fresh sample $z\sim\mathcal{N}(0,I)$
produces something plausible. A plain autoencoder has no such guarantee---its
latent space may be full of holes that decode to nonsense.

\subsection{Restricted Boltzmann Machines}

An \keyterm{RBM} is a bipartite energy-based model over visible units $v$ and
hidden units $h$, with no connections inside a layer:
\begin{equation}
  E(v,h) = -v^{\!\top}Wh - b^{\!\top}v - c^{\!\top}h,
  \qquad P(v,h)\propto e^{-E(v,h)} .
  \label{eq:dl-rbm}
\end{equation}
\where{$v$ is the vector of visible units, $h$ the hidden units, $W$ the
weight matrix coupling them, and $b,c$ the visible and hidden biases. Lower
energy means higher probability; the missing normalising constant is what makes
exact training intractable.}
The bipartite restriction makes $P(h\mid v)$ and $P(v\mid h)$ factorise into
independent Bernoullis, which is what makes sampling tractable. Exact maximum
likelihood is not tractable, because it requires the partition function;
\keyterm{contrastive divergence} approximates the
gradient\textsuperscript{\cite[][]{Kriesel2007NeuralNetworks}} by contrasting a
data-clamped phase with a short Gibbs-sampled model phase. RBMs powered the
greedy layer-wise pre-training that made deep belief nets trainable in
2006\textsuperscript{\cite{hintonFastLearningAlgorithm2006}}---historically
decisive, and now superseded by better initialisation and normalisation.

\begin{checkpoint}
For each of GAN, VAE, and autoencoder, state (a) what you can sample from it,
(b) whether it gives you a likelihood, and (c) the failure mode you would check
first. Which one would you choose to detect anomalous sensor readings, and why?
\end{checkpoint}

% ===========================================================================
\section{Research Directions Beyond Backpropagation}
% ===========================================================================

\subsection{The Weight Transport Problem and Feedback Alignment}

Backpropagation sends error signals backwards using the \emph{transpose} of the
forward weights, $W^{\!\top}$\textsuperscript{\cite[][p.~3]{fischerVorlesungGrundlagenNeuronale2024}}.
The backward path must therefore know every forward weight exactly---the
\keyterm{weight transport problem}, for which no biological mechanism is known.

Lillicrap and colleagues showed that replacing $W^{\!\top}$ with a \emph{fixed
random} matrix $B$ still supports
learning\textsuperscript{\cite{lillicrapRandomSynapticFeedback2016}}: the
forward weights gradually align with $B$, so the network shapes itself to make
the random feedback useful. \keyterm{Direct feedback alignment} goes further,
delivering the error from the output layer to every hidden layer through
separate random matrices, bypassing the chain
entirely\textsuperscript{\cite{noklandDirectFeedbackAlignment2016}}.

\subsection{Synthetic Gradients}

In standard training no layer may update until the full forward and backward
passes finish---\keyterm{update locking}. \keyterm{Decoupled neural interfaces}
attach a small predictor $M_i$ to each layer that \emph{estimates} the incoming
gradient, letting the layer update
immediately\textsuperscript{\cite{jaderbergDecoupledNeuralInterfaces2017}}. Each
$M_i$ is then trained against the true gradient when it eventually arrives.

\begin{figure}[htbp]
\centering
\begin{tikzpicture}[
  node distance=6mm and 10mm,
  lay/.style={draw=AlmanachTeal,fill=AlmanachMist,rounded corners=1mm,
    minimum width=17mm,minimum height=6.5mm,font=\sffamily\scriptsize},
  dni/.style={draw=AlmanachOchre,fill=AlmanachOchre!22,rounded corners=1mm,
    minimum width=15mm,minimum height=6mm,font=\sffamily\scriptsize},
  arr/.style={-{Stealth[length=1.8mm]},thick,draw=AlmanachInk!55},
  darr/.style={-{Stealth[length=1.8mm]},dashed,draw=AlmanachRed},
  lbl/.style={font=\sffamily\scriptsize,color=AlmanachInk!70}]
  \node[lay] (L1) {Layer 1};
  \node[lay,right=of L1] (L2) {Layer 2};
  \node[lay,right=of L2] (L3) {Layer 3};
  \node[lay,right=of L3] (loss) {Loss};
  \node[dni,below=of L1] (M1) {DNI $M_1$};
  \node[dni,below=of L2] (M2) {DNI $M_2$};
  \draw[arr] (L1) -- (L2) node[midway,above,lbl] {$h_1$};
  \draw[arr] (L2) -- (L3) node[midway,above,lbl] {$h_2$};
  \draw[arr] (L3) -- (loss);
  \draw[arr] (L1) -- (M1); \draw[arr] (L2) -- (M2);
  \draw[darr] (M1) -- (L1)
    node[midway,right,lbl,color=AlmanachRed] {$\tilde\delta_1$};
  \draw[darr] (M2) -- (L2)
    node[midway,right,lbl,color=AlmanachRed] {$\tilde\delta_2$};
  \node[lbl,below=3mm of M1,align=center,text width=42mm]
    {each layer updates on a \emph{predicted} gradient
     without waiting for the loss};
\end{tikzpicture}
\caption{Decoupled neural interfaces. The dashed arrows are synthetic
gradients supplied locally, which removes update locking and allows layers to
train asynchronously---at the cost of learning from an estimate rather than the
true error signal.}
\label{fig:dl-synthetic-gradients}
\end{figure}

\subsection{Capsules and Spiking Networks}

\keyterm{Capsule networks} address a real limitation of pooling: it buys
invariance by discarding pose\textsuperscript{\cite{sabourDynamicRoutingCapsules2017}}.
A capsule is a group of neurons whose activity \emph{vector} encodes both
presence (its length) and instantiation parameters such as pose and orientation
(its direction). An iterative \keyterm{dynamic routing} procedure sends each
lower capsule's output to whichever higher capsule its prediction agrees with,
yielding \emph{equivariance}---the representation changes predictably as the
viewpoint changes. The cost is quadratic routing, and the approach has not
scaled to large datasets.

\keyterm{Spiking neural networks} communicate through discrete events rather
than continuous activations\textsuperscript{\cite{maassNetworksSpikingNeurons1997}}.
A leaky integrate-and-fire unit accumulates input in a membrane potential and
emits a spike when it crosses threshold. Sparse events map naturally onto
neuromorphic hardware at far lower power than dense GPU arithmetic. The
obstacle is training: the spike function is not differentiable, so
\keyterm{surrogate gradients} substitute a smooth approximation in the backward
pass only.

\begin{engineernote}
Everything in this section is scientifically interesting and, as of the date on
this book's title page, none of it displaces backpropagation in production.
Report it as an open research direction, not as an available option. The honest
summary is: these methods answer ``could learning work differently?'' with yes,
and ``should you use them tomorrow?'' with no.
\end{engineernote}

% ===========================================================================
\section{Deep Dive: The Economics of Training and Inference}
% ===========================================================================

\subsection{Pre-Norm Residual Blocks}

Modern Transformers normalise \emph{before} each sub-layer rather than after:
\[
  x' = x + \mathrm{Attention}\bigl(\mathrm{Norm}(x)\bigr),
  \qquad
  y = x' + \mathrm{MLP}\bigl(\mathrm{Norm}(x')\bigr).
\]
The residual stream then runs from input to output without passing through any
normalisation, so the identity path of \cref{eq:dl-residual} stays exactly
intact. This is what allows very deep stacks to train without a carefully tuned
warmup.

\subsection{Why Generation Is Slower Than Training}

Self-attention processes all training positions in parallel, but autoregressive
generation is irreducibly sequential: token $t+1$ cannot be computed before
token $t$. A \keyterm{key--value cache} stores the keys and values already
computed so each new token attends to history without recomputing it---trading
memory for latency.

\begin{workedexample}[Sizing a key--value cache]
Take a model with $L = 30$ layers, $H_{kv} = 3$ key/value heads, head dimension
$d_h = 64$, in \texttt{float16} ($2$ bytes). Per token the cache holds keys
\emph{and} values in every layer:
\[
  2 \times L \times H_{kv} \times d_h \times 2\,\text{bytes}
  = 2 \times 30 \times 3 \times 64 \times 2
  = 23{,}040\ \text{bytes} \approx 22.5\ \text{KiB per token.}
\]
For a $4{,}096$-token context:
\[
  4096 \times 23{,}040\ \text{bytes} \approx 94.4\ \text{MB.}
\]
With a batch of $32$ concurrent requests that is about $3.0\,\text{GB}$ of cache
alone---on top of the weights. \emph{This is why serving throughput is usually
memory-bound rather than compute-bound}, and why grouped-query attention (here
$3$ key/value heads shared by $9$ query heads) exists at all: it divides this
number by the sharing factor. The same architecture is examined end to end in
\cref{ch:practical-llm-engineering}.
\end{workedexample}

\subsection{Scaling Laws, and Their Correction}

Model quality follows smooth power laws in parameters, data, and
compute\textsuperscript{\cite{kaplanScalingLawsNeural2020}}. That result was
widely read as ``make the model bigger''. Hoffmann and colleagues later showed
that the models of that era were substantially
\emph{under-trained}: for a fixed compute budget, parameters and training tokens
should grow in roughly equal
proportion\textsuperscript{\cite{hoffmannTrainingComputeOptimal2022}}.

\begin{engineernote}
This is a case worth remembering for its methodology, not only its result. Both
papers were careful and empirical; the second overturned the practical reading
of the first by varying a factor the first had held nearly fixed. Scaling laws
are fitted curves over a studied range, not laws of nature. Quote them with
their range, their date, and their architecture, or do not quote them.
\end{engineernote}

\begin{table}[htbp]
\centering
\caption{Efficiency techniques and the cost each one actually reduces. No entry
improves everything, and every entry introduces an error or a constraint that
must be measured rather than assumed.}
\label{tab:dl-efficiency}
\begin{tblr}{
  width=\textwidth,
  colspec={Q[l,0.19\textwidth] X[l] X[l]},
  row{1}={font=\bfseries,bg=AlmanachMist},
  hlines,vlines}
Technique & Reduces & Introduces \\
Quantisation & Weight memory, bandwidth &
  Rounding error, worst on outlier activation channels \\
Distillation & Parameters, latency &
  A student ceiling set by the teacher's own errors \\
Pruning / sparsity & Parameters, sometimes compute &
  Irregular access patterns that hardware may not exploit \\
LoRA\textsuperscript{\cite{huLoRALowRank2021}} & Trainable parameters,
  optimiser memory & A rank limit on what adaptation can express \\
Activation checkpointing & Training memory &
  Roughly $30\,\%$ extra compute from recomputation \\
FlashAttention\textsuperscript{\cite{daoFlashAttentionFastMemory2022}} &
  Memory traffic, peak memory &
  Nothing in accuracy---it is exact---but it is kernel- and
  hardware-specific \\
KV cache & Repeated computation & Memory linear in context $\times$ batch \\
\end{tblr}
\end{table}

\begin{practicallab}[Measure the trade-off instead of assuming it]
\textbf{Goal.} Replace the sentence ``quantisation is basically free'' with a
number you measured yourself.

\textbf{Setup.} Any small instruct model you can run locally---the checkpoint
used in \cref{ch:practical-llm-engineering} is a good choice---plus a fixed
prompt set of at least $50$ items with reference answers.

\textbf{Protocol.}
\begin{enumerate}
  \item Record the environment: exact model revision, library versions,
        hardware, and random seed. Without this the rest is anecdote.
  \item Establish the baseline in full precision. Measure quality on your
        prompt set, time to first token, decode tokens per second, and peak
        memory.
  \item Repeat with 8-bit and 4-bit quantisation, changing \emph{one} variable
        at a time.
  \item Plot quality against peak memory. Mark the configuration where quality
        first drops below your acceptance threshold.
  \item Run each configuration three times and report the spread, not just the
        mean.
\end{enumerate}

\textbf{Deliverable.} A one-page dossier with a table of four measurements per
configuration, the variance across repeats, and one sentence naming the
configuration you would deploy and the evidence for it.

\textbf{The point.} Parameter count is not a deployment metric. Quality,
latency, throughput, peak memory, energy, and operational complexity each move
differently, and only measurement tells you which one binds in your setting.
\end{practicallab}

\begin{checkpoint}
Using the key--value cache calculation, work out the maximum batch size that
fits in $16\,\text{GB}$ of memory at a $8{,}192$-token context, ignoring
weights. Then state which single architectural change would double it.
\end{checkpoint}

\chapterreview{Represent}{Differentiate}{Optimise}{Scale}

% ===========================================================================
\section{Further Reading}
% ===========================================================================

\begin{v3topic}{Primary sources for the mechanisms in this chapter}
\begin{itemize}
  \item \fullcite{rumelhartLearningRepresentationsBackpropagating1986} ---
    the paper that made multi-layer training practical.
  \item \fullcite{vaswaniAttentionAllYou2023} --- the Transformer:
    self-attention, multi-head attention, positional encoding.
  \item \fullcite{heDeepResidualLearning2016} --- residual learning, the change
    that removed the depth barrier.
  \item \fullcite{hochreiterLongShortTermMemory1997} --- the LSTM, and the
    additive memory path.
  \item \fullcite{kingmaAdamMethodStochastic2017} and
    \fullcite{loshchilovDecoupledWeightDecay2019} --- the optimiser you will
    actually use, and the correction that made it right.
  \item \fullcite{hoffmannTrainingComputeOptimal2022} --- compute-optimal
    scaling, and a model case of a careful result revising an earlier one.
\end{itemize}
\end{v3topic}

\begin{v3topic}{Textbook treatments}
\begin{itemize}
  \item \fullcite{fleuretLittleBookDeep} --- compact and unusually well judged
    on what to leave out.
  \item \fullcite{Kriesel2007NeuralNetworks} --- biological foundations,
    perceptrons, and RBMs with historical context.
  \item \fullcite{robertsPrinciplesDeepLearning2022} --- a theoretical account
    of initialisation, depth, and generalisation.
  \item \fullcite{zhangUnderstandingDeepLearning2021} --- why the classical
    generalisation story is incomplete.
\end{itemize}
\end{v3topic}
